\documentclass[11pt,reqno]{amsart}
\usepackage[localtheorems]{raeez-math-template}

% ================================================================
% Packages
% ================================================================
\usepackage{array}

% ================================================================
% Theorem environments
% ================================================================
\newtheorem{theorem}{Theorem}[section]
\newtheorem{proposition}[theorem]{Proposition}
\newtheorem{lemma}[theorem]{Lemma}
\newtheorem{corollary}[theorem]{Corollary}
\newtheorem{conjecture}[theorem]{Conjecture}
\theoremstyle{definition}
\newtheorem{definition}[theorem]{Definition}
\newtheorem{construction}[theorem]{Construction}
\newtheorem{example}[theorem]{Example}
\newtheorem{computation}[theorem]{Computation}
\theoremstyle{remark}
\newtheorem{remark}[theorem]{Remark}
\newtheorem{warning}[theorem]{Warning}

% ================================================================
% Macros
% ================================================================
\newcommand{\cA}{\mathcal{A}}
\newcommand{\cE}{\mathcal{E}}
\newcommand{\cM}{\mathcal{M}}
\newcommand{\cW}{\mathcal{W}}
\newcommand{\Mbar}{\overline{\mathcal{M}}}
\newcommand{\fg}{\mathfrak{g}}
\newcommand{\fh}{\mathfrak{h}}
\newcommand{\fn}{\mathfrak{n}}
\newcommand{\fsl}{\mathfrak{sl}}
\newcommand{\gmod}{\mathfrak{g}^{\mathrm{mod}}_\cA}
\newcommand{\gEone}{\mathfrak{g}^{E_1}_\cA}
\newcommand{\Eone}{E_1}
\newcommand{\Einf}{E_\infty}
\newcommand{\Ash}{\cA^{\mathrm{sh}}}
\newcommand{\Defcyc}{\mathrm{Def}_{\mathrm{cyc}}}
\newcommand{\Ran}{\mathrm{Ran}}
\newcommand{\MC}{\mathrm{MC}}
\newcommand{\Sym}{\mathrm{Sym}}
\newcommand{\Conf}{\mathrm{Conf}}
\newcommand{\End}{\mathrm{End}}
\newcommand{\Ad}{\mathrm{Ad}}
\newcommand{\Res}{\mathrm{Res}}
\newcommand{\Aut}{\mathrm{Aut}}
\newcommand{\Sh}{\mathrm{Sh}}
\newcommand{\Vir}{\mathrm{Vir}}
\newcommand{\Heis}{\mathrm{Heis}}
\newcommand{\av}{\mathrm{av}}
\newcommand{\dfib}{d_{\mathrm{fib}}}
\newcommand{\CW}{\mathrm{CW}}
\providecommand{\reg}{\mathrm{reg}}
\providecommand{\Z}{\mathbb{Z}}

\DeclareMathOperator{\tr}{tr}
\DeclareMathOperator{\rk}{rk}
\DeclareMathOperator{\depth}{depth}
\DeclareMathOperator{\obs}{obs}
\DeclareMathOperator{\im}{Im}

\numberwithin{equation}{section}

% ================================================================
\begin{document}

\title[Chiral Chern--Weil Theory]
{Chiral Chern--Weil Theory}

\author{Raeez Lorgat}
\address{Perimeter Institute for Theoretical Physics,
 Waterloo, ON N2L 2Y5, Canada}
\email{rlorgat@perimeterinstitute.ca}

\date{}

\begin{abstract}
Classical Chern--Weil theory extracts characteristic classes from
the curvature of a principal bundle using invariant polynomials.
We develop a chiral analogue in the setting of modular Koszul
duality for vertex algebras. The ordered bar complex
$B^{\mathrm{ord}}(\cA)$ over the Ran space of a curve~$X$
plays the role of a chiral classifying space; the Maurer--Cartan
element $\Theta_\cA^{\Eone}$ plays the role of a universal
connection, with binary component the collision kernel
$K_\cA^{\mathrm{coll}}(z)$.  The homotopy-coinvariant comparison
produces $q_n((\Theta_\cA^{\Eone})_{0,n})$, and specified normalized
functionals extract scalar shadow data by moment-residue, Sugawara,
and composite-sector rules.  A represented classical $r$-matrix
appears under $H_{\mathrm{CYBE}}^{\mathrm{rep}}$. As the curve~$X$
varies over $\Mbar_g$, the fiberwise bar differential has an Arakelov
curvature term.  For $g\geq2$, the resulting Theorem~D data consist
of $\operatorname{Obs}^{\mathrm{def}}_g(\cA)$, the perfect object
$\mathfrak O_g^K(\cA)=\kappa(\cA)\lambda_{-1}(\mathbb E_g)$, its
signed Hodge character, and the stable-graph functional.
The stable genus-one component begins with
$\operatorname{Obs}^{\mathrm{def}}_{1,1}(\cA)$.
The degree-by-degree scalar coordinates
\[
S_n(\cA)=\tau_{\cA,n}^{\mathrm{res}}
q_n((\Theta_\cA^{E_1})_{0,n})
\]
form the shadow tower, with
$(S_2,S_3,S_4)=(\kappa,\mathfrak C,\mathfrak Q)$ on the corresponding
scalar channels. The four standard depth types (G/L/C/M) are
read from the tuple $(S_3,S_4,\text{termination})$, not from the
single downstream discriminant $\Delta=8\kappa S_4$.
For multi-weight algebras, the scalar formula conditionally acquires a
cross-channel correction
$\delta F_g^{\mathrm{cross}}$, which we compute explicitly
for $\cW_3$ at genus~$2$:
$\delta F_2(\cW_3) = (c{+}204)/(16c)$.
\end{abstract}

\subjclass[2020]{17B69, 14H10, 18M70, 81T40, 57R20}

\keywords{Chiral algebras, Chern--Weil theory, bar construction,
Koszul duality, moduli of curves, $r$-matrix, vertex algebras}

\maketitle

\setcounter{tocdepth}{2}
\tableofcontents


% ================================================================
% Section 1: Introduction
% ================================================================
\section{Introduction}\label{sec:intro}

\subsection{The deficiency}
The bar complex of a chiral algebra~$\cA$ on a curve~$X$
varies as $X$ moves in the moduli space $\Mbar_g$.
At genus~$0$, the bar differential squares to zero.
At genus~$1$, it does not: the fiberwise differential acquires
fiber Arakelov curvature, and its pushforward to the base produces
the Hodge-class obstruction with coefficient~$\kappa$.

This curvature carries $\Eone$ operadic indices.  For an affine
Kac--Moody algebra $V_k(\fg)$, the binary collision kernel
$K_{k,\fg}^{\mathrm{coll}}(z)=k\,\Omega_{\mathrm{tr}}/z$ is valued in
$\fg\otimes\fg$.  Its homotopy coinvariant
$q_2(K_{k,\fg}^{\mathrm{coll}})$ remains operator-valued, and the
normalized residue trace supplies the scalar coordinate.  For the
Virasoro algebra at central charge~$c$, the collision kernel
$K_{\mathrm{Vir}_c}^{\mathrm{coll}}(z)=(c/2)/z^3+2T/z$ contains both
vacuum and field-dependent terms; the moment-residue functional gives
$\tau_{\mathrm{Vir}_c,2}^{\mathrm{res}}q_2
(K_{\mathrm{Vir}_c}^{\mathrm{coll}})=c/2$.

The question is: what mechanism extracts the scalar shadow data from
$K_\cA^{\mathrm{coll}}$?
The answer has a precise classical analogue. In classical
Chern--Weil theory, the curvature of a principal $G$-bundle is
$\fg$-valued, and a $G$-invariant polynomial kills the
internal indices to produce a scalar characteristic class.
The chiral analogue has two successive maps.  The averaging comparison
$q_n\colon\operatorname{End}_\cA(n)\to
\operatorname{End}_\cA(n)_{h\Sigma_n}$ produces an equivariant operator
class.  A specified normalized functional
$\tau_{\cA,n}^{\mathrm{res}}$ then produces its scalar coordinate.

\subsection{Results}
The main contributions are:

\begin{enumerate}[(1)]
\item A systematic parallel between classical Chern--Weil theory
 and the modular Koszul duality programme
 (Sections~\ref{sec:classical}--\ref{sec:chiral-cw}).

\item Explicit computation of $q_2$ and
 $\tau_{\cA,2}^{\mathrm{res}}$ on the collision kernels of the
 standard families: Heisenberg, Virasoro,
 affine Kac--Moody (all types $A$--$G$), principal $\cW_N$,
 and free-field systems (Section~\ref{sec:averaging}).

\item Identification of the shadow tower
 $(\kappa, \mathfrak{C}, \mathfrak{Q}, \ldots)$ as the
 sequence of chiral characteristic classes, with the G/L/C/M
 classification as the classification of chiral classifying
 spaces by the type of their characteristic ring
 (Section~\ref{sec:shadow}).

\item Derivation and explicit computation of the multi-weight
 cross-channel correction $\delta F_g^{\mathrm{cross}}$,
 including the first closed-form result
 $\delta F_2(\cW_3) = (c{+}204)/(16c)$
 (Section~\ref{sec:multi-weight}).
\end{enumerate}


% ================================================================
% Section 2: Classical Chern--Weil
% ================================================================
\section{Classical Chern--Weil theory}\label{sec:classical}

\subsection{Connections and curvature}

Let $G$ be a Lie group with Lie algebra~$\fg$, and let
$P \to M$ be a principal $G$-bundle over a smooth manifold~$M$.
The bundle may be topologically non-trivial, but any single
local trivialization sees only~$\fg$. To detect the global
topology from local data, one needs a device that is defined
locally but transforms coherently: a connection.

\begin{definition}\label{def:connection}
A \emph{connection} on~$P$ is a $\fg$-valued $1$-form
$\nabla \in \Omega^1(P;\fg)$ satisfying:
\begin{enumerate}[(i)]
\item $R_g^*\nabla = \Ad(g^{-1})\nabla$ for all $g \in G$
 (equivariance);
\item $\nabla(\xi_P) = \xi$ for all $\xi \in \fg$
 (normalization along fibers),
\end{enumerate}
where $\xi_P$ is the fundamental vector field of~$\xi$
and $R_g$ is the right $G$-action on~$P$.
In a local trivialization over $U \subset M$, a connection
is represented by $A \in \Omega^1(U;\fg)$.
\end{definition}

\begin{definition}\label{def:curvature}
The \emph{curvature} of a connection~$\nabla$ is the
$\fg$-valued $2$-form
\begin{equation}\label{eq:curvature-def}
F_\nabla = d\nabla + \tfrac{1}{2}[\nabla,\nabla]
 \;\in\; \Omega^2(M;\Ad(P)).
\end{equation}
In local coordinates: $F = dA + A \wedge A$.
\end{definition}

The curvature measures the failure of parallel transport around
infinitesimal loops to be trivial. It is $\fg$-valued: it
carries internal indices from the Lie algebra.

\begin{lemma}[Bianchi identity]\label{lem:bianchi}
$d_\nabla F_\nabla = 0$.
\end{lemma}

\begin{proof}
Compute $d_\nabla F = d_\nabla(d\nabla +
\frac{1}{2}[\nabla,\nabla]) = d^2\nabla +
\frac{1}{2}[d_\nabla\nabla,\nabla] -
\frac{1}{2}[\nabla,d_\nabla\nabla]$. Using $d^2 = 0$,
$d_\nabla\nabla = F$, and the Jacobi identity for~$\fg$,
all terms cancel.
\end{proof}

\subsection{The Chern--Weil homomorphism}

\begin{definition}\label{def:inv-poly}
Let $\Sym^k(\fg^*)^G$ denote the space of $G$-invariant
symmetric $k$-linear forms on~$\fg$. For
$P \in \Sym^k(\fg^*)^G$, the \emph{Chern--Weil form} is
\begin{equation}\label{eq:cw-form}
\CW(P) = P(\underbrace{F_\nabla, \ldots, F_\nabla}_k)
 \;\in\; \Omega^{2k}(M).
\end{equation}
\end{definition}

The form $\CW(P)$ is scalar-valued: the invariant polynomial
contracts all $\fg$-indices. Its construction requires the
curvature to carry indices (otherwise there is nothing to
contract) and the polynomial to be invariant (otherwise the
result would depend on the trivialization).

\begin{theorem}[Chern--Weil]\label{thm:cw}
\begin{enumerate}[\upshape(i)]
\item $\CW(P)$ is closed: $d[\CW(P)] = 0$.
\item The cohomology class $[\CW(P)] \in H^{2k}(M;\mathbb{R})$
 is independent of the connection~$\nabla$.
\item The resulting map
 $\CW\colon \Sym^*(\fg^*)^G \to H^*(M;\mathbb{R})$
 is a ring homomorphism.
\end{enumerate}
\end{theorem}

\begin{proof}
(i)~The Bianchi identity $d_\nabla F = 0$ and the
$G$-invariance of~$P$ imply
$d[P(F^k)] = k\,P(d_\nabla F, F^{k-1}) = 0$.

\smallskip\noindent
(ii)~Let $\nabla_0$, $\nabla_1$ be connections with curvatures
$F_0$, $F_1$. Set $\alpha = \nabla_1 - \nabla_0 \in
\Omega^1(M;\Ad(P))$ and interpolate:
$\nabla_t = \nabla_0 + t\alpha$,
$F_t = F_0 + t\,d_{\nabla_0}\alpha + t^2\alpha\wedge\alpha$.
Then
\[
P(F_1^k) - P(F_0^k) = d\!\int_0^1 k\,P(\alpha, F_t^{k-1})\,dt.
\]
The Chern--Simons transgression
$\mathrm{TP} = \int_0^1 k\,P(\alpha, F_t^{k-1})\,dt$
exhibits the difference as exact.

\smallskip\noindent
(iii)~Multiplicativity of the wedge product.
\end{proof}

\begin{example}[Abelian: $G = U(1)$]\label{ex:u1}
$\fg = \mathbb{R}$, $F = dA$ is already scalar. The unique
invariant polynomial (up to scale) is the identity. The first
Chern class $c_1 = \frac{1}{2\pi i}[F] \in H^2(M;\mathbb{Z})$
is the sole characteristic class.
The classifying space $BU(1) = \mathbb{CP}^\infty$ has
$H^*(BU(1)) = \mathbb{Z}[c_1]$.
\end{example}

\begin{example}[$G = SU(2)$]\label{ex:su2}
$\fg = \mathfrak{su}(2)$ is traceless, so $c_1 = 0$.
The first non-trivial invariant polynomial is
$P(X,Y) = -\frac{1}{8\pi^2}\tr(XY)$, giving the
second Chern class $c_2 = -\frac{1}{8\pi^2}[\tr(F\wedge F)]$.
On~$S^4$, $\int_{S^4}c_2 \in \mathbb{Z}$ is the instanton
number.
\end{example}

\begin{remark}[Structure of the invariant ring]\label{rem:type}
For $G = GL_n(\mathbb{C})$, the invariant polynomials are
generated by $P_k(X) = \tr(X^k)$, $k = 1,\ldots,n$, producing
Chern classes $c_1,\ldots,c_n$. The number of independent
generators of $\Sym^*(\fg^*)^G$ determines the \emph{type} of
the classifying space~$BG$.
\end{remark}


% ================================================================
% Section 3: Chiral setup
% ================================================================
\section{Chiral algebras, the bar complex, and the
collision kernel}\label{sec:chiral}

Classical Chern--Weil theory requires three ingredients: a
connection, a curvature, and an invariant polynomial. Each
has a chiral counterpart, and each requires construction.
The connection is the Maurer--Cartan element of the $\Eone$
convolution algebra; its binary twisting component is the collision
kernel.  The homotopy-coinvariant projection supplies an equivariant
operator class, and the normalized residue trace plays the role of an
invariant polynomial.  A representation satisfying
$H_{\mathrm{CYBE}}^{\mathrm{rep}}$ identifies the represented collision
kernel with a classical $r$-matrix.

\subsection{Chiral algebras and OPE}

Let~$X$ be a smooth algebraic curve. A chiral algebra~$\cA$
on~$X$ is a factorization algebra encoding operator product
expansions. In local coordinates, the data consists of
strong generators $\varphi_1,\ldots,\varphi_r$ of conformal
weights $h_1,\ldots,h_r$, with operator product expansion
\begin{equation}\label{eq:ope}
\varphi_i(z)\,\varphi_j(w) \;\sim\;
\sum_{p=1}^{h_i + h_j}
\frac{c^p_{ij}(w)}{(z-w)^p}
\end{equation}
where $c^p_{ij}(w)$ are OPE coefficients, possibly depending
on the remaining fields.

\subsection{The ordered bar complex}

\begin{definition}\label{def:bar}
The \emph{ordered bar complex} of~$\cA$ is
\[
B^{\mathrm{ord}}(\cA) = T^c(s^{-1}\bar\cA)
\]
where $\bar\cA = \ker(\varepsilon)$ is the augmentation ideal
of~$\cA$, $s^{-1}$ is the desuspension ($|s^{-1}v| = |v| - 1$),
and $T^c$ is the cofree conilpotent coalgebra with the
deconcatenation coproduct. Grading is cohomological
($|d| = +1$).
\end{definition}

\begin{warning}
Three distinct errors are common in the bar-complex formula.
First, using~$\cA$ instead of~$\bar\cA$ (the augmentation ideal
is essential).
Second, using the suspension~$s$ instead of
desuspension~$s^{-1}$ (bar is down).
Third, using the tensor algebra~$T$ instead of the cofree
coalgebra~$T^c$ (the deconcatenation coproduct is the
coalgebra structure).
\end{warning}

On a curve~$X$, the bar complex is a factorization coalgebra
on the Ran space $\Ran(X) = \bigsqcup_{n \ge 1}\Conf_n(X)$.
The bar differential uses the logarithmic propagator
$\eta_{12} = d\log(z_1 - z_2)$ at genus~$0$.

\subsection{The collision kernel from the OPE}

The binary bar differential encodes the collision of two generators.
The propagator $d\log(z_1-z_2)$ absorbs one pole from the OPE and gives
the collision kernel.

\begin{definition}\label{def:r-matrix}
The \emph{collision kernel} of~$\cA$ is the binary component
of the $\Eone$ Maurer--Cartan element:
$\Theta^{\Eone}_\cA \in \MC(\gEone)$:
\[
K_\cA^{\mathrm{coll}}(z)
\;=\; (\Theta^{\Eone}_\cA)_{0,2}
\;\in\; \operatorname{End}_\cA(2)\otimes\cO(*\Delta).
\]
It is obtained from the OPE data by $d\log$-absorption; its pole order
equals $p_{\max}(\mathrm{OPE})-1$.  Given a representation
$\rho\colon\operatorname{End}_\cA(2)\to\operatorname{End}(V^{\otimes2})$
and $H_{\mathrm{CYBE}}^{\mathrm{rep}}$, the operator
$\rho(K_\cA^{\mathrm{coll}})$ is a classical $r$-matrix.
\end{definition}

We now compute the collision kernel for every standard family.

\begin{computation}[Heisenberg at level~$k$]
\label{comp:r-heis}
Generator~$J$, weight $h = 1$. OPE:
\[
J(z)J(w) = \frac{k}{(z-w)^2}.
\]
Pole order~$2$. After $d\log$-absorption:
\begin{equation}\label{eq:r-heis}
\boxed{\;K_{\Heis_k}^{\mathrm{coll}}(z) = \frac{k}{z}.\;}
\end{equation}
This is a simple pole.  At $k=0$ the kernel is zero.
\end{computation}

\begin{computation}[Virasoro at central charge~$c$]
\label{comp:r-vir}
Generator~$T$, weight $h = 2$. OPE:
\[
T(z)T(w) = \frac{c/2}{(z-w)^4}
+ \frac{2T(w)}{(z-w)^2} + \frac{\partial T(w)}{z-w}.
\]
Pole order~$4$. After $d\log$-absorption, only the singular
part of the OPE contributes to the collision residue; the
regular $\partial T$ term drops (its $d\log$ coefficient
vanishes identically). Thus:
\begin{equation}\label{eq:r-vir}
\boxed{\;K_{\Vir_c}^{\mathrm{coll}}(z)
\;=\; \frac{c/2}{z^3} \;+\; \frac{2T}{z}.\;}
\end{equation}
Cubic pole, \emph{not} quartic. The leading term
$(c/2)/z^3$ is a $c$-number; the subleading term
$2T/z$ is field-dependent. The regular term $\partial T$
from the OPE is absent in $K_{\Vir_c}^{\mathrm{coll}}$: it is $d\log$-absorbed,
or equivalently it is a total derivative on the configuration
space and contributes $0$ to the residue pairing.
At central charge $c=0$, the vacuum cubic term vanishes and the
field-dependent piece $2T/z$ remains.
\end{computation}

\begin{computation}[Affine Kac--Moody $V_k(\fg)$]
\label{comp:r-km}
Generators $J^a$, $a = 1,\ldots,\dim(\fg)$, all weight
$h = 1$. Let $(\ ,\ )$ be the normalized invariant form
(long roots squared length~$2$), $\{t^a\}$ an orthonormal
basis. OPE:
\[
J^a(z)J^b(w) = \frac{k\,\delta^{ab}}{(z-w)^2}
 + \frac{f^{ab}{}_c\,J^c(w)}{z-w}.
\]
After $d\log$-absorption:
\begin{equation}\label{eq:r-km}
\boxed{\;K_{k,\fg}^{\mathrm{coll}}(z) = \frac{k\,\Omega}{z},\quad
 \Omega = \sum_a t^a \otimes t^a \in \fg\otimes\fg.\;}
\end{equation}
This is the \emph{trace-form convention}. The level~$k$
survives $d\log$-absorption. At $k=0$ the trace-form collision kernel
is zero.  In the represented \emph{Kazhdan--Lusztig (KZ) convention},
under $H_{\mathrm{CYBE}}^{\mathrm{rep}}$,
$r_{k,\fg}^{\mathrm{KZ}}(z)=\Omega_{\mathrm{KZ}}/((k+h^\vee)z)$
where $\Omega_{\mathrm{KZ}}$ is the Killing-normalized Casimir;
at $k = 0$ this gives $\Omega_{\mathrm{KZ}}/(h^\vee z) \neq 0$
for non-abelian $\fg$ (the Lie bracket persists). The two
conventions are \emph{not equal} as rational functions of $k$:
they are related by a level-dependent rescaling
$\Omega_{\mathrm{tr}} = (2 h^\vee)^{-1}\Omega_{\mathrm{KZ}}$
together with $k \mapsto k(k+h^\vee)/(2h^\vee)$, or
equivalently by absorbing the $(k+h^\vee)^{-1}$ factor into
the KZ flat connection. The equation
$k\,\Omega_{\mathrm{tr}} = \Omega_{\mathrm{KZ}}/(k+h^\vee)$
should be read as a structural \emph{change of basis} between
conventions, not as an identity of rational functions in $k$:
at $k=0$ the LHS vanishes while the RHS is finite and non-zero.
We adopt the trace convention throughout this note; every affine
collision kernel in this section carries its level factor explicitly.
\end{computation}

\begin{computation}[Principal $\cW_N$ algebra]
\label{comp:r-wn}
Generators $W^{(s)}$ at weights $s = 2,3,\ldots,N$
($N-1$ generators). The self-OPE of the spin-$s$ primary
in the Fateev--Lukyanov normalization has leading singularity
\[
W^{(s)}(z)\,W^{(s)}(w) = \frac{c/s}{(z-w)^{2s}} + \cdots
\]
After $d\log$-absorption, the vacuum leading term is
\begin{equation}\label{eq:r-wn}
K_{ss}^{\mathrm{vac}}(z) = \frac{c/s}{z^{2s-1}} + \cdots
\end{equation}
The cross-OPE $W^{(s_1)}(z)\,W^{(s_2)}(w)$ for $s_1 \neq s_2$
has leading coefficient that is field-dependent, not a $c$-number.
\end{computation}

\begin{computation}[Fermionic $bc$ system at weight~$\lambda$]
\label{comp:r-bc}
Generators: $b$ (weight~$\lambda$) and $c$
(weight~$1-\lambda$), both fermionic.
Central charge $c_{bc}(\lambda) = 1 - 3(2\lambda-1)^2$.
The singular OPEs:
\[
b(z)\,c(w) \sim \frac{1}{z-w},\qquad
b(z)\,b(w) \sim 0,\qquad
c(z)\,c(w) \sim 0.
\]
The only singular OPE has a simple pole (pole order~$1$).
After $d\log$-absorption the cross-contraction has pole order~$0$.
It does not contribute to the scalar collision residue:
\begin{equation}\label{eq:r-bc}
K_{bc}^{\mathrm{coll}}(z) = 0.
\end{equation}
The self-OPEs $b(z)b(w)$ and $c(z)c(w)$ are regular
(fermionic statistics forbid self-contraction at
generic~$\lambda$), so $K_{bb}^{\mathrm{coll}}=K_{cc}^{\mathrm{coll}}=0$.

At $\lambda = 2$: the reparametrization ghost system of the
bosonic string, with $c_{bc}(2) = -26$.
At $\lambda = 1/2$: a single Dirac fermion, $c_{bc}(1/2) = 1$.
\end{computation}

\begin{computation}[Bosonic $\beta\gamma$ system at
weight~$\lambda$]\label{comp:r-bg}
Generators: $\beta$ (weight~$\lambda$) and $\gamma$
(weight~$1-\lambda$), both bosonic.
Central charge
$c_{\beta\gamma}(\lambda) = 2(6\lambda^2 - 6\lambda + 1)
= -c_{bc}(\lambda)$.
OPE:
\[
\beta(z)\,\gamma(w) \sim \frac{1}{z-w}.
\]
After $d\log$-absorption the cross-contraction has pole order~$0$,
so $K_{\beta\gamma}^{\mathrm{coll}}(z)=0$ in the collision-residue
shadow. The modular characteristic comes from the stress tensor and
composite sector.

Complementarity check: $c_{\beta\gamma}(\lambda) +
c_{bc}(\lambda) = 0$ for all~$\lambda$.
At $\lambda = 1$: $c_{\beta\gamma}(1) = 2$,
$c_{bc}(1) = -2$, sum~$= 0$.
At $\lambda = 2$: $c_{\beta\gamma}(2) = 26$,
$c_{bc}(2) = -26$, sum~$= 0$ (bosonic string ghost
cancellation).
\end{computation}


% ================================================================
% Section 4: The averaging map
% ================================================================
\section{The averaging map and the modular
characteristic}\label{sec:averaging}

\subsection{Definition of the averaging map}

\begin{definition}[Averaging comparison]\label{def:av}
An \emph{averaging datum} is an $R$-twisted descent datum from the
ordered ($\Eone$) convolution Lie algebra to the modular symmetric
convolution Lie algebra. With such a datum fixed, the arity-$n$
comparison is the quotient to $\Sigma_n$-coinvariants:
\begin{equation}\label{eq:av-def}
\av\colon \gEone \longrightarrow \gmod.
\end{equation}
For finite-dimensional aritywise calculations over characteristic
zero, the Reynolds operator
\[
\mathrm{Reyn}(\phi)(g,n)=\frac1{n!}\sum_{\sigma\in\Sigma_n}
\sigma\cdot\phi(g,n)
\]
identifies invariants with coinvariants after a choice of splitting;
it is not a canonical dg Lie section.
\end{definition}

\begin{theorem}[Averaging on the descent locus]\label{thm:av-properties}
On the descent locus, the comparison map satisfies:
\begin{enumerate}[\upshape(i)]
\item $\av$ is a dg~Lie morphism.
\item $\av(\Theta^{\Eone}_\cA) = \Theta_\cA$: the ordered
 MC element maps to the symmetric MC element.
\item In arity~$n$, the image is the equivariant operator class
 $q_n((\Theta_\cA^{\Eone})_{g,n})$; its scalar coordinate is defined
 after choosing $\tau_{\cA,n}^{\mathrm{res}}$.
\end{enumerate}
\end{theorem}

At the level of MC elements, the ordered MC element has degree-wise
components
$(K_\cA^{\mathrm{coll}},(\Theta_\cA^{\Eone})_{0,3},
(\Theta_\cA^{\Eone})_{0,4},\ldots)$.  The Maurer--Cartan equation gives
their operadic boundary relations.  Under
$H_{\mathrm{CYBE}}^{\mathrm{rep}}$, the represented binary component
satisfies the classical Yang--Baxter equation; the corresponding
arity-three comparison carries the represented associator relation.

The modular MC element $\Theta_\cA=\av(\Theta^{\Eone}_\cA)$ has
arity-$n$ component $q_n((\Theta_\cA^{\Eone})_{g,n})$ in the homotopy
coinvariants.

\subsection{The $\Sigma_2$ action on a collision kernel}
\label{subsec:sigma2}

Before computing~$\kappa$ for specific families, we work
out what the averaging map does to a generic collision kernel at
degree~$2$, step by step.

At degree~$2$, the $\Eone$ convolution algebra consists of
elements on the ordered configuration space
$\Conf_2(X) = \{(z_1,z_2) \in X^2 : z_1 \neq z_2\}$.
Let $z = z_1 - z_2$ be the relative coordinate.
After choosing a representation on the space~$V$ of strong generators,
the collision kernel is a section of $\End(V^{\otimes 2})$ over
$\Conf_2(X)$.
For generators $\varphi_i$, $\varphi_j$, write
\begin{equation}\label{eq:generic-r}
K^{\mathrm{coll}}(z) = \sum_{i,j} K_{ij}(z)\cdot
e_i \otimes e_j
\;\in\;
\End(V^{\otimes 2})\otimes\mathbb{C}((z)),
\end{equation}
where $\{e_i\}$ is a basis for~$V$ and each~$K_{ij}(z)$ is a
Laurent series in~$z$ (possibly with field-dependent
coefficients).

The group $\Sigma_2$ acts on $\Conf_2(X)$ by swapping:
$(z_1,z_2) \mapsto (z_2,z_1)$, equivalently $z \mapsto -z$.
Simultaneously, $\Sigma_2$ acts on $\End(V)^{\otimes 2}$ by
swapping tensor factors: $e_i \otimes e_j \mapsto
e_j \otimes e_i$. The combined action on the collision kernel is:
\begin{equation}\label{eq:sigma2-action}
(\sigma \cdot K^{\mathrm{coll}})(z) = \sum_{i,j} K_{ij}(-z)\cdot
e_j \otimes e_i.
\end{equation}

Under the characteristic-zero Reynolds identification, the
$\Sigma_2$-coinvariant class has representative
\begin{equation}\label{eq:sigma2-coinvariant}
q_2(K^{\mathrm{coll}})(z)
= \tfrac{1}{2}\bigl(K^{\mathrm{coll}}(z)
+ \sigma\cdot K^{\mathrm{coll}}(z)\bigr)
= \tfrac{1}{2}\sum_{i,j}
\bigl(K_{ij}(z) + K_{ji}(-z)\bigr)\cdot
e_i \otimes e_j.
\end{equation}

\subsection{From function to scalar: the cyclic trace}
\label{subsec:cyclic-trace}

The class~\eqref{eq:sigma2-coinvariant} is an
$\End(V^{\otimes2})$-valued function of~$z$.  Residue extraction and
the chosen cyclic trace combine into
$\tau_{\cA,2}^{\mathrm{res}}$.

\medskip\noindent
\emph{Residue extraction.}
The modular convolution algebra at degree~$2$ pairs
with the propagator $\eta_{12} = d\log(z_1 - z_2)$, which
extracts the residue at $z = 0$. The relevant coefficient is
\begin{equation}\label{eq:residue-step}
\Res_{z=0}\bigl[q_2(K^{\mathrm{coll}})(z)\bigr]
= \tfrac{1}{2}\sum_{i,j}
\bigl(\Res_{z=0}\,K_{ij}(z) + \Res_{z=0}\,K_{ji}(-z)\bigr)
\cdot e_i\otimes e_j.
\end{equation}
For an even-order pole $K_{ij}(z) = c/z^{2m}$:
$\Res_{z=0}(c/z^{2m}) = 0$ and
$\Res_{z=0}(c/(-z)^{2m}) = 0$.
For an odd-order pole $K_{ij}(z) = c/z^{2m+1}$:
$\Res_{z=0}(c/z^{2m+1}) = 0$ unless $m = 0$
(simple pole), and
$\Res_{z=0}(c/(-z)^{2m+1})
= (-1)^{2m+1}\Res_{z=0}(c/z^{2m+1})
= -\Res_{z=0}(c/z^{2m+1})$.

For a simple-pole collision kernel $K_{ij}(z) = c_{ij}/z$:
\[
\Res_{z=0}\,c_{ij}/z + \Res_{z=0}\,c_{ji}/(-z)
= c_{ij} - c_{ji}.
\]
This is the \emph{antisymmetric} part of the coefficient
matrix. For a symmetric coefficient ($c_{ij} = c_{ji}$,
as for $k\,\delta^{ab}$ in the Heisenberg/KM double pole),
the two terms cancel!

The resolution: for higher-weight generators, the relevant
coefficient is not the simple residue but the
\emph{moment residue}, as we now explain.

\medskip\noindent
\emph{The moment and the cyclic trace.}
The cyclic trace on $\End(V^{\otimes 2})$ incorporates the
conformal weights of the generators. For a generator
$\varphi_i$ of weight~$h_i$, the OPE pairing carries a
weight factor $z^{2h_i - 2}$ from the conformal Ward identity:
the two-point function
$\langle\varphi_i(z)\varphi_i(0)\rangle \propto z^{-2h_i}$
contributes a volume factor $z^{2h_i}$, and the propagator
$d\log z$ contributes~$z^{-2}$, giving a net factor
$z^{2h_i - 2}$.

The cyclic trace therefore evaluates the
\emph{$(2h_i{-}2)$-th moment} of the residue:
\begin{equation}\label{eq:cyclic-trace-moment}
\tau_{\cA,2}^{\mathrm{res}}q_2(K_\cA^{\mathrm{coll}})
= \sum_i \Res_{z=0}\bigl[z^{2h_i - 2}\cdot
 K_{ii}^{\mathrm{vac}}(z)\bigr]
+ (\text{Sugawara shift}).
\end{equation}

For a single generator of weight~$h$ with vacuum collision kernel
$K^{\mathrm{vac}}(z) = c_h/z^{2h-1} + \cdots$
(leading vacuum pole), the moment residue is:
\begin{equation}\label{eq:moment-residue}
\Res_{z=0}\bigl[z^{2h-2}\cdot c_h/z^{2h-1}\bigr]
= \Res_{z=0}\bigl[c_h/z\bigr] = c_h.
\end{equation}
The moment factor $z^{2h-2}$ exactly cancels the
excess pole order, reducing any leading pole to a simple pole
whose residue is the vacuum OPE coefficient~$c_h$.

\medskip\noindent
\emph{Field-dependent terms.}
The full collision kernel contains field-dependent terms in addition
to the vacuum ($c$-number) part. For Virasoro:
$K_{\mathrm{Vir}_c}^{\mathrm{coll}}(z) = (c/2)/z^3 + 2T/z$, where $2T/z$
is field-dependent.

The vacuum projection
$K^{\mathrm{vac}}(z) = \langle 0|K^{\mathrm{coll}}(z)|0\rangle$ retains
the $c$-number part. The cyclic trace annihilates
field-dependent terms because they have zero vacuum expectation
value: $\langle 0|T|0\rangle = 0$. More precisely, the
trace on $\End(V^{\otimes 2})$ contracts the two $V$-slots
against the vacuum state, and any term proportional to a
non-identity field gives zero after contraction.

For Virasoro at $h = 2$:
\[
\Res_{z=0}\bigl[z^2\cdot K_{\mathrm{Vir}_c}^{\mathrm{coll}}(z)\bigr]
= \Res_{z=0}\bigl[z^2\cdot\bigl((c/2)/z^3 + 2T/z\bigr)\bigr]
= \Res_{z=0}\bigl[(c/2)/z + 2Tz\bigr]
= c/2 + 0.
\]
The field-dependent term $2T$ is regular
at $z = 0$ after the moment shift, so their residue vanishes.
This is a general mechanism: the moment factor $z^{2h-2}$
shifts the vacuum pole to a simple pole and simultaneously
makes all field-dependent terms regular.

\medskip
For families whose scalar shadow is purely diagonal, the extraction
rule is
\begin{equation}\label{eq:kappa-full}
\kappa(\cA)
= \sum_i \kappa_i + \kappa_{\mathrm{sp}},
\end{equation}
where $\kappa_i = \Res_{z=0}[z^{2h_i-2}\cdot
K_{ii}^{\mathrm{vac}}(z)]$ is the per-channel diagonal
contribution and $\kappa_{\mathrm{sp}}$ is the Sugawara shift
from simple-pole self-contraction (nonzero only for
non-abelian same-weight generators).
First-order systems and non-principal reductions also contain
composite-sector contributions, so \eqref{eq:kappa-full} is an
extraction rule on the diagonal standard families, not a universal
trace-residue theorem.

\subsection{Explicit trace: $\mathfrak{sl}_2$ at level~$k$}
\label{subsec:sl2-explicit}

We carry out the full computation for $\fg = \mathfrak{sl}_2$,
$\dim = 3$, $h^\vee = 2$. The standard basis is
\[
e = \begin{pmatrix} 0 & 1 \\ 0 & 0 \end{pmatrix},\quad
f = \begin{pmatrix} 0 & 0 \\ 1 & 0 \end{pmatrix},\quad
h = \begin{pmatrix} 1 & 0 \\ 0 & -1 \end{pmatrix},
\]
with bracket $[e,f] = h$, $[h,e] = 2e$, $[h,f] = -2f$.

The normalized invariant form (long roots squared
length~$2$) is $(x,y) = \frac{1}{2h^\vee}\tr_{\mathrm{ad}}
(\mathrm{ad}(x)\,\mathrm{ad}(y))
= \frac{1}{4}\tr_{\mathrm{ad}}
(\mathrm{ad}(x)\,\mathrm{ad}(y))$.
Compute: $(e,f) = 1$, $(h,h) = 2$, $(e,e) = (f,f) = 0$.
The form is indefinite; no real orthonormal basis exists.

The Casimir is $\Omega = \sum_a t^a \otimes t_a$ where
$\{t^a\}$ and $\{t_a\}$ are dual bases for~$(\ ,\ )$.
With $t^1 = e$, $t^2 = f$, $t^3 = h/2$ and
$t_1 = f$, $t_2 = e$, $t_3 = h/2$
(so that $(t^a, t_b) = \delta^a_b$):
\begin{equation}\label{eq:sl2-casimir}
\Omega = e \otimes f + f \otimes e
+ \tfrac{1}{2}\,h \otimes h.
\end{equation}
Check: $\tr(\Omega) = (e,f) + (f,e) + \frac{1}{2}(h,h)
= 1 + 1 + 1 = 3 = \dim(\mathfrak{sl}_2)$.

The collision kernel is
$K_{k,\mathfrak{sl}_2}^{\mathrm{coll}}(z)
=k\,\Omega_{\mathrm{tr}}/z$.

\medskip\noindent
\emph{Double-pole channel.}
Each generator~$J^a$ has weight $h = 1$; the moment factor is
$z^0 = 1$. The double-pole OPE
$J^a(z)\,J_a(w) \sim k\,(t^a,t_a)/(z-w)^2 = k/(z-w)^2$
gives collision-kernel coefficient $k/z$ per dual pair. Three dual pairs,
total double-pole trace:
\[
\tr\bigl(\dfib^{\,2}\big|_{\mathrm{dp}}\bigr)
= 3k\cdot\omega_1.
\]
Normalizing by $2h^\vee = 4$:
$\kappa_{\mathrm{dp}} = 3k/(2h^\vee) = 3k/4$.

\medskip\noindent
\emph{Simple-pole channel.}
Use the normalized invariant form with dual bases
$\{x_i\}$ and $\{x^i\}$, so
\[
\Omega=\sum_i x_i\otimes x^i
=e\otimes f+f\otimes e+\tfrac12h\otimes h .
\]
The adjoint Casimir identity is
\[
\sum_i[x_i,[x^i,y]]=2h^\vee y .
\]
For $\mathfrak{sl}_2$, $h^\vee=2$; with
$x_i=e,f,h$ and $x^i=f,e,h/2$, the identity gives
$\sum_i[x_i,[x^i,e]]=4e$. This is the simple-pole
self-contraction responsible for the Sugawara shift.

The Sugawara trace:
\[
\kappa_{\mathrm{sp}} = \frac{\dim(\fg)}{2} = \frac{3}{2}.
\]

\medskip\noindent
\emph{Total.}
\[
\kappa(V_k(\mathfrak{sl}_2))
= \kappa_{\mathrm{dp}} + \kappa_{\mathrm{sp}}
= \frac{3k}{4} + \frac{3}{2}
= \frac{3(k+2)}{4}
= \frac{\dim(\fg)(k+h^\vee)}{2h^\vee}.\quad\checkmark
\]
At $k = 1$: $\kappa = 9/4$. At $k = -2$ (critical):
$\kappa = 0$. At $k = 0$:
$\kappa = 3/2 = \dim(\mathfrak{sl}_2)/2$
(the Sugawara contribution from the three generators).

\subsection{Explicit trace: $\mathfrak{sl}_3$ at level~$k$}
\label{subsec:sl3-explicit}

For $\fg = \mathfrak{sl}_3$: $\dim = 8$, $h^\vee = 3$.
The Lie algebra has the Chevalley basis
$\{e_1,e_2,e_3,f_1,f_2,f_3,h_1,h_2\}$ (three positive roots,
three negative roots, two Cartan generators), giving eight
generators $J^a$ for $a = 1,\ldots,8$.

The collision kernel is
$K_{k,\mathfrak{sl}_3}^{\mathrm{coll}}(z)
=k\,\Omega_{\mathrm{tr}}/z$ with
$\Omega_{\mathrm{tr}} = \sum_{a=1}^8
t^a \otimes t^a$.

\medskip\noindent
\emph{Double-pole channel.}
Each of the $8$ generators has weight~$1$, moment factor
$z^0 = 1$, and per-channel residue
$\kappa_a^{\mathrm{dp}} = k$. Summing and normalizing:
\[
\kappa_{\mathrm{dp}} = \frac{8k}{2h^\vee} = \frac{8k}{6}
= \frac{4k}{3}.
\]

\medskip\noindent
\emph{Simple-pole channel.}
The adjoint Casimir eigenvalue is
$C_2^{\mathrm{ad}} = 2h^\vee = 6$. The Sugawara trace:
\[
\kappa_{\mathrm{sp}} = \frac{\dim(\fg)}{2} = \frac{8}{2} = 4.
\]

\medskip\noindent
\emph{Total.}
\[
\kappa(V_k(\mathfrak{sl}_3)) = \frac{4k}{3} + 4
= \frac{4k + 12}{3} = \frac{4(k+3)}{3}
= \frac{\dim(\fg)(k+h^\vee)}{2h^\vee}.\quad\checkmark
\]
At $k = 1$: $\kappa = 16/3$. At $k = -3$ (critical):
$\kappa = 0$.

\subsection{Computation: $bc$ and $\beta\gamma$ systems}
\label{subsec:av-bc-bg}

The $bc$ and $\beta\gamma$ systems have a distinct-weight
structure ($h_b = \lambda \neq 1-\lambda = h_c$ for
$\lambda \neq 1/2$) and their only singular OPE is a
cross-OPE.

\medskip\noindent
\emph{The $bc$ system.}
Generators $b$ (weight~$\lambda$) and $c$
(weight~$1{-}\lambda$).
The cross-contraction has pole order~$0$ after $d\log$-absorption:
\[
\Res_{z=0}[K_{bc}^{\mathrm{coll}}(z)] = 0.
\]
The collision residue has \emph{no residue}. The kappa
the diagonal collision-kernel extraction gives zero.

How does $\kappa$ arise? For the $bc$ system, the stress
tensor is
\[
T_{bc} = (\partial b)\,c\cdot\lambda + b\,(\partial c)
\cdot(1-\lambda),
\]
a composite field of weight~$2$. The curvature
$\dfib^{\,2} = \kappa\cdot\omega_g$ arises not from the
diagonal collision kernel directly
but from the Sugawara-type construction: the composite stress
tensor generates a Virasoro subalgebra with central charge
$c_{bc}(\lambda) = 1 - 3(2\lambda-1)^2$.

The modular characteristic of the full $bc$ vertex algebra is
\begin{equation}\label{eq:kappa-bc}
\kappa(bc_\lambda) = \frac{c_{bc}(\lambda)}{2}
= \frac{1 - 3(2\lambda-1)^2}{2}.
\end{equation}
This arises from the Virasoro subalgebra generated by
$T_{bc}$: the bar complex of the $bc$ system, as a
factorization coalgebra, has fiberwise curvature controlled by
the stress tensor, and $\kappa = c/2$ when the curvature is
entirely generated by the Virasoro sector.

Boundary checks:
$\lambda = 1/2 \Rightarrow c_{bc} = 1$,
$\kappa = 1/2$ (single Dirac fermion).
$\lambda = 2 \Rightarrow c_{bc} = -26$,
$\kappa = -13$ (string reparametrization ghost).

\medskip\noindent
\emph{The $\beta\gamma$ system.}
By complementarity $c_{\beta\gamma} + c_{bc} = 0$:
\begin{equation}\label{eq:kappa-bg}
\kappa(\beta\gamma_\lambda) = \frac{c_{\beta\gamma}(\lambda)}{2}
= \frac{2(6\lambda^2-6\lambda+1)}{2}
= 6\lambda^2-6\lambda+1.
\end{equation}
Checks: $\lambda = 1/2 \Rightarrow \kappa = -1/2$;
$\lambda = 2 \Rightarrow \kappa = 13$;
$\kappa(bc_\lambda) + \kappa(\beta\gamma_\lambda)
= (1-3(2\lambda-1)^2)/2 + (6\lambda^2-6\lambda+1) = 0$
for all~$\lambda$ (Koszul complementarity, conductor~$K = 0$).

\begin{remark}\label{rem:bc-r-matrix}
The $bc$ and $\beta\gamma$ systems illustrate a subtlety: when
the only singular OPE is a \emph{cross}-OPE (not a self-OPE),
the diagonal collision kernel has zero residue.  The scalar~$\kappa$
arises from the Sugawara stress tensor of the
composite system: the bar complex detects the conformal anomaly
through the Virasoro sector, regardless of whether the
generators have non-trivial self-OPEs. For free-field systems
whose OPE is purely a cross-OPE, $\kappa = c/2$ exactly, where
$c$ is the central charge of the stress tensor. This is
consistent with the general formula~\eqref{eq:kappa-full}:
the ``Sugawara shift'' subsumes the entire contribution when
the diagonal collision kernel is zero.
\end{remark}

\subsection{Degree-$2$: diagonal extraction rule}
\label{subsec:kappa-general}

\begin{proposition}\label{prop:kappa-formula}
For diagonal standard families, the scalar degree-$2$ contribution is
extracted by the weight-dependent moment residue. For a single
generator~$\varphi_i$ of weight~$h_i$, the per-channel contribution is
\begin{equation}\label{eq:kappa-per-channel}
\kappa_i = \Res_{z=0}\bigl[z^{2h_i - 2}\cdot
 K_{ii}^{\mathrm{vac}}(z)\bigr].
\end{equation}
For abelian and Virasoro-type families,
$\kappa(\cA)=\sum_i\kappa_i$. For non-abelian affine
Kac--Moody in the trace-form convention,
$\kappa(\cA)=\sum_i\kappa_i+\dim(\fg)/2$.
\end{proposition}

\subsection{Computation: Heisenberg}\label{subsec:av-heis}

$K_{\Heis_k}^{\mathrm{coll}}(z)=k/z$, $h=1$.
\[
\kappa_J = \Res_{z=0}\bigl[z^{2\cdot 1 - 2}\cdot k/z\bigr]
= \Res_{z=0}[k/z] = k.
\]

\begin{equation}\label{eq:kappa-heis}
\boxed{\;\kappa(H_k) = k.\;}
\end{equation}

Boundary checks: $k = 0 \Rightarrow \kappa = 0$ (uncurved);
$k = 1 \Rightarrow \kappa = 1$.

\subsection{Computation: Virasoro}\label{subsec:av-vir}

$K_{\Vir_c}^{\mathrm{vac}}(z)=(c/2)/z^3$, $h=2$.
\[
\kappa_T = \Res_{z=0}\bigl[z^{2\cdot 2 - 2}\cdot
 (c/2)/z^3\bigr]
= \Res_{z=0}\bigl[z^2\cdot(c/2)/z^3\bigr]
= \Res_{z=0}\bigl[(c/2)/z\bigr] = c/2.
\]

\begin{equation}\label{eq:kappa-vir}
\boxed{\;\kappa(\Vir_c) = c/2.\;}
\end{equation}

Boundary checks: $c = 0 \Rightarrow \kappa = 0$;
$c = 13 \Rightarrow \kappa = 13/2$ (self-dual point);
$c = 26 \Rightarrow \kappa = 13$.

\subsection{Computation: $\cW_N$}\label{subsec:av-wn}

Generators $W^{(s)}$ at weights $s = 2,\ldots,N$ with
per-channel vacuum collision kernel
$K_{ss}^{\mathrm{vac}}(z)=(c/s)/z^{2s-1}+\cdots$.

For each generator:
\[
\kappa_s = \Res_{z=0}\bigl[z^{2s-2}\cdot
 (c/s)/z^{2s-1}\bigr]
= \Res_{z=0}\bigl[(c/s)/z\bigr] = c/s.
\]

Sum over generators:
\begin{equation}\label{eq:kappa-wn}
\boxed{\;\kappa(\cW_N) = \sum_{s=2}^{N}\frac{c}{s}
 = c\cdot(H_N - 1),\quad
 H_N = \sum_{j=1}^{N}\frac{1}{j}.\;}
\end{equation}

Verification at $N=2$ ($\cW_2 = \Vir$):
$\kappa = c\cdot(3/2 - 1) = c/2$.

Verification at $N=3$:
$\kappa = c\cdot(1/2 + 1/3) = 5c/6$.

Verification at $N=4$:
$\kappa = c\cdot(1/2 + 1/3 + 1/4) = 13c/12$.

\begin{warning}\label{warn:harmonic}
$H_{N-1} \neq H_N - 1$. At $N = 2$: $H_1 = 1$ but
$H_2 - 1 = 1/2$. The formula uses $H_N - 1$ (subtract the
integer~$1$ from the $N$-th harmonic number), not $H_{N-1}$
(the harmonic number at~$N-1$).
\end{warning}

\subsection{Computation: affine Kac--Moody and the
Sugawara shift}\label{subsec:av-km}

For $V_k(\fg)$ with $\fg$ simple, the trace-form collision kernel is
$K_{k,\fg}^{\mathrm{coll}}(z)=k\,\Omega_{\mathrm{tr}}/z$.
The functional $\tau_{k,\fg,2}^{\mathrm{res}}$ evaluates its
double-pole channel and gives $\kappa_{\mathrm{dp}}$.  A one-loop
self-contraction in the simple-pole OPE channel supplies the Sugawara
term and completes~$\kappa$.

The proof proceeds in three steps.

\medskip\noindent
\emph{Double-pole channel.}
The diagonal OPE $J^a(z)J^a(w) \sim k/(z-w)^2$ contributes
to the fiberwise curvature $\dfib^{\,2}$ on a genus-$1$ curve:
\begin{equation}\label{eq:dp-channel}
\dfib^{\,2}\big|_{\mathrm{dp}}(J^a \boxtimes J^b)
= (-2\pi i)\cdot k\,(J^a,J^b)\cdot\omega_1,
\end{equation}
where $\omega_1 = c_1(\mathbb{E})$ is the Hodge class.
Taking the trace over the normalized form $\sum_a(J^a,J_a)
= \dim(\fg)$:
\begin{equation}\label{eq:dp-trace}
\tr\bigl(\dfib^{\,2}\big|_{\mathrm{dp}}\bigr)
= k\cdot\dim(\fg)\cdot\omega_1.
\end{equation}

\medskip\noindent
\emph{Simple-pole channel (one-loop self-contraction).}
The cross-OPE $J^a(z)J^b(w) \sim f^{ab}{}_c\,J^c(w)/(z-w)$
feeds back into a second OPE contraction: $J^c$ is
contracted with itself through a second propagator loop,
forming a one-loop Feynman diagram.
\begin{equation}\label{eq:sp-channel}
\dfib^{\,2}\big|_{\mathrm{sp}}(J^a \boxtimes J^b)
= (-2\pi i)\cdot\tfrac{1}{2}\sum_{c,d}
 f^{ac}{}_d\,f^{bc}{}_d\cdot\omega_1.
\end{equation}
The factor~$\frac{1}{2}$ is the symmetry factor
($|\Aut| = 2$ of the one-loop diagram). The contraction
$\sum_{c,d}f^{ac}{}_d\,f^{bc}{}_d$ is the adjoint
quadratic Casimir eigenvalue:
\begin{equation}\label{eq:casimir-adj}
\sum_{c,d}f^{ac}{}_d\,f^{bc}{}_d
= C_2^{\mathrm{ad}}\cdot(J^a,J^b) = 2h^\vee\cdot(J^a,J^b).
\end{equation}
This is the Killing form identity: $K^{ab} = 2h^\vee\,
\delta^{ab}$ in an orthonormal basis. Taking the trace:
\begin{equation}\label{eq:sp-trace}
\tr\bigl(\dfib^{\,2}\big|_{\mathrm{sp}}\bigr)
= \tfrac{1}{2}\cdot 2h^\vee\cdot\dim(\fg)\cdot\omega_1
= h^\vee\cdot\dim(\fg)\cdot\omega_1.
\end{equation}

\medskip\noindent
\emph{Combining the channels.}
On a single diagonal component, the two channels give
per-component curvature $k + h^\vee$. Expressing the
curvature as a multiple of the Casimir
$\mathsf{C} = \sum_a J^a\otimes J_a$:
\[
m_0 = \frac{k+h^\vee}{2h^\vee}\cdot\mathsf{C},
\]
where $1/(2h^\vee)$ normalizes against $C_2^{\mathrm{ad}}$.
Then $\kappa = \tr(m_0)$:
\begin{align}
\kappa &= \frac{k+h^\vee}{2h^\vee}\cdot\tr(\mathsf{C})
= \frac{k+h^\vee}{2h^\vee}\cdot\sum_a(J^a,J_a)
\notag\\
&= \frac{(k+h^\vee)\cdot\dim(\fg)}{2h^\vee}.
\label{eq:kappa-km-derivation}
\end{align}

The decomposition:
\begin{equation}\label{eq:dp-sp-decomposition}
\kappa = \underbrace{\frac{k\cdot\dim(\fg)}{2h^\vee}}_{%
 \kappa_{\mathrm{dp}}\text{ (tree-level collision-kernel trace)}}
\;+\;
\underbrace{\frac{\dim(\fg)}{2}}_{%
 \kappa_{\mathrm{sp}}\text{ (one-loop, Sugawara shift)}}.
\end{equation}

\begin{equation}\label{eq:kappa-km}
\boxed{\;\kappa\bigl(V_k(\fg)\bigr)
= \frac{\dim(\fg)\,(k + h^\vee)}{2h^\vee}.\;}
\end{equation}

\begin{example}[Consistency checks]\label{ex:km-checks}
\begin{center}
\begin{tabular}{@{}llllll@{}}
\toprule
$\fg$ & $\dim$ & $h^\vee$ & $\kappa$ & $k=1$ &
 $k=-h^\vee$ \\
\midrule
$\mathfrak{sl}_2$ & $3$ & $2$ & $3(k{+}2)/4$ & $9/4$ & $0$ \\
$\mathfrak{sl}_3$ & $8$ & $3$ & $4(k{+}3)/3$ & $16/3$ & $0$ \\
$E_6$ & $78$ & $12$ & $13(k{+}12)/4$ & $169/4$ & $0$ \\
$E_7$ & $133$ & $18$ & $133(k{+}18)/36$ & $2527/36$ & $0$ \\
$E_8$ & $248$ & $30$ & $62(k{+}30)/15$ & $1922/15$ & $0$ \\
$G_2$ & $14$ & $4$ & $7(k{+}4)/4$ & $35/4$ & $0$ \\
$F_4$ & $52$ & $9$ & $26(k{+}9)/9$ & $260/9$ & $0$ \\
\bottomrule
\end{tabular}
\end{center}
At $k = -h^\vee$ (critical level), $\kappa = 0$: the
uncurved fixed point of Feigin--Frenkel.
At $k = 0$, $\kappa = \dim(\fg)/2$: the Sugawara shift
alone.
\end{example}

\subsection{When the Sugawara shift vanishes}
\label{subsec:sugawara-vanishing}

\begin{proposition}\label{prop:sugawara-vanishing}
The Sugawara shift $\kappa_{\mathrm{sp}}$ is non-zero if and
only if the algebra has multiple generators of the same
conformal weight with non-trivial cross-OPE at the
simple-pole level. In the scalar-shift-free cases,
$\kappa=\tau_{\cA,2}^{\mathrm{res}}q_2(K_\cA^{\mathrm{coll}})$.
\end{proposition}

Three cases:

\begin{enumerate}[(i)]
\item \emph{Single-generator algebras}
 (Heisenberg, Virasoro): no cross-OPE,
 $\kappa_{\mathrm{sp}} = 0$.

\item \emph{Distinct-weight multi-generator algebras}
 ($\cW_N$): the $c$-number contributions come only from
 diagonal self-OPEs; the cross-OPEs between generators of
 different weights have field-dependent leading coefficients
 that do not produce closed loops.
 $\kappa_{\mathrm{sp}} = 0$.

\item \emph{Same-weight multi-generator algebras}
 (Kac--Moody): all $\dim(\fg)$ generators have weight~$1$,
 the structure constants~$f^{ab}{}_c$ create the one-loop
 self-contraction, and
 $\kappa_{\mathrm{sp}} = \dim(\fg)/2 \neq 0$.
\end{enumerate}

\subsection{Summary of collision kernels and $\kappa$-values}
\label{subsec:summary-table}

\begin{center}
\begin{tabular}{@{}llll@{}}
\toprule
Family & $K^{\mathrm{coll}}(z)$ (vacuum leading) & $\kappa$ &
$\kappa_{\mathrm{sp}}$ \\
\midrule
$H_k$ & $k/z$ & $k$ & $0$ \\
$\Vir_c$ & $(c/2)/z^3$ & $c/2$ & $0$ \\
$V_k(\fg)$ & $k\,\Omega_{\mathrm{tr}}/z$ &
 $\dim(\fg)(k{+}h^\vee)/(2h^\vee)$ & $\dim(\fg)/2$ \\
$\cW_N$ & $(c/s)/z^{2s-1}$ per spin & $c(H_N{-}1)$ & $0$ \\
$bc_\lambda$ & $0$ in collision residue &
 $(1{-}3(2\lambda{-}1)^2)/2$ & all from Sug. \\
$\beta\gamma_\lambda$ & $0$ in collision residue &
 $6\lambda^2{-}6\lambda{+}1$ & all from Sug. \\
\bottomrule
\end{tabular}
\end{center}


% ================================================================
% Section 5: The chiral curvature
% ================================================================
\section{The chiral curvature}\label{sec:chiral-cw}

The averaging comparison sends the ordered collision kernel to the
operator class $q_2(K_\cA^{\mathrm{coll}})$.  On abelian and scalar
families the moment-residue functional gives
$\tau_{\cA,2}^{\mathrm{res}}q_2(K_\cA^{\mathrm{coll}})=\kappa$.
For non-abelian affine Kac--Moody algebras the same functional gives
$\kappa_{\mathrm{cl}}$, and the Sugawara term completes the modular
characteristic:
$\kappa=\kappa_{\mathrm{cl}}+\dim(\fg)/2$.
The question it forces is: what role does~$\kappa$ play
geometrically? The answer: $\kappa$ is the proportionality
constant of a curvature. The bar differential, which squares
to zero on~$\mathbb{P}^1$, fails to square to zero on
higher-genus curves; the failure is proportional to~$\kappa$.
This section derives that formula from the propagator on a
genus-$g$ curve.

\subsection{Propagators at higher genus}

On~$\mathbb{P}^1$, the propagator
$\eta_{12} = d\log(z_1 - z_2)$ is single-valued and the bar
differential satisfies $d_{\mathrm{bar}}^2 = 0$. On a
genus-$g$ Riemann surface $\Sigma_g$ with $g \ge 1$, the
function $z_1 - z_2$ is not globally defined. The
single-valued (Arakelov-normalized) propagator is
\begin{equation}\label{eq:propagator-g}
\eta^{(g)}_{ij} = \Bigl[\partial_{z_i}\!\log E(z_i,z_j)
+ \pi\sum_{\alpha,\beta=1}^{g}
\omega_\alpha(z_i)\,(\im\,\Omega)^{-1}_{\alpha\beta}\,
\im\!\Bigl(\int_{z_0}^{z_j}\!\omega_\beta\Bigr)
\Bigr](dz_i - dz_j),
\end{equation}
where $E(z,w)$ is the prime form, $\omega_1,\ldots,\omega_g$
are the normalized abelian differentials, and~$\Omega$ is the
$g\times g$ period matrix.

At genus~$1$ ($\Sigma_g = E_\tau$ with period~$\tau$):
\begin{equation}\label{eq:propagator-1}
\eta^{(1)}_{12}
=
\left[
\partial_u\log\vartheta_1(u\mid\tau)
+2\pi i\,\frac{\im u}{\im\tau}
\right]_{u=z_1-z_2}
(dz_1-dz_2),
\end{equation}
where $\vartheta_1$ is the odd Jacobi theta function. The
real-analytic coefficient multiplies the same $(1,0)$ differential as
the theta quotient; its $+2\pi i$ shift under $u\mapsto u+\tau$
cancels the theta quotient's $-2\pi i$ shift, giving a
double-periodic Arakelov-normalized propagator.

\subsection{The curvature formula}

\begin{theorem}[Fiber curvature and base obstruction]\label{thm:curvature}
The fiberwise bar differential on $B(\cA)$ over~$\Sigma_g$
satisfies
\begin{equation}\label{eq:curvature}
\dfib^{\,2} = \kappa(\cA)\cdot\omega_g^{\mathrm{Ar}},
\end{equation}
where $\omega_g^{\mathrm{Ar}}$ is the Arakelov $(1,1)$-form on the
fiber curve. The GRR/Deligne-pairing pushforward, under $H_D^K$,
produces
\[
 \mathfrak O_g^K(\cA)=\kappa(\cA)\lambda_{-1}(\mathbb E_g),
 \qquad
 \operatorname{ch}_g(\mathfrak O_g^K(\cA))
 =(-1)^g\kappa(\cA)c_g(\mathbb E_g).
\]
For $g\geq2$, the native class
$\operatorname{Obs}^{\mathrm{def}}_g(\cA)\in H^2(\Def_g(\cA))$
enters this formula through the comparison package
$H_D^{\mathrm{tr}}$.
\end{theorem}

At genus~$1$:
$\omega_1 = \frac{i}{2\im(\tau)}\,dz\wedge d\bar z$.
At genus~$g \ge 2$:
$\omega_g$ involves the full period matrix
$(\im\,\Omega)^{-1}$, which is a $g \times g$ matrix
(not a scalar).

The curvature factorizes:
\[
\dfib^{\,2} = \underbrace{\kappa(\cA)}_{\text{algebra}}
\;\cdot\;
\underbrace{\omega_g}_{\text{geometry}}.
\]
The first factor is intrinsic to the chiral algebra
(a quasi-isomorphism invariant, independent of the curve).
The second factor is intrinsic to the complex structure of
the curve (independent of the algebra). This factorization
is a structural feature of the chiral theory with no direct
classical Chern--Weil analogue.

\subsection{Cardy--Arakelov test for \texorpdfstring{$\kappa$}{kappa}}
\label{subsec:cardy-arakelov}

The curvature formula~\eqref{eq:curvature} is algebraic: $\kappa(\cA)$
is extracted by $\tau_{\cA,2}^{\mathrm{res}}q_2$ from the collision
kernel, and $\omega_g$
is a geometric form on $\Sigma_g$. The Cardy--Arakelov identity
exhibits the same $\kappa$ as the Arakelov-normalised coefficient of
the coincident-limit torus two-point function of the stress tensor,
thereby providing an independent geometric derivation of $\kappa$
from $\langle T(z)T(w)\rangle_{\tau}$.

\begin{conjecture}[Cardy--Arakelov test for the Virasoro lane]
\label{thm:cardy-arakelov-kappa}
Let $\cA$ be a chirally Koszul algebra whose scalar invariant is the
Virasoro lane $\kappa=c/2$, and let $\Sigma_1 = E_\tau$ be an elliptic curve
with Arakelov $(1,1)$-form
$\omega_1 = \tfrac{i}{2\im(\tau)}\,dz\wedge d\bar z$
\textup{(}\textup{}: $\int_{E_\tau}\omega_1 = 1$\textup{)}. Let
$\langle T(z)T(w)\rangle_\tau^{\reg}$ denote the torus two-point
function regularised by subtracting the $c/2\cdot(z-w)^{-4}$
holomorphic short-distance singularity and the Sugawara
self-energy shift~$\Delta T_{\mathrm{Sug}}$. Then the modular
Virasoro-lane scalar satisfies
\begin{equation}\label{eq:cardy-arakelov}
\frac{c}{2}
\;=\;
\lim_{w\to z}
\frac{\langle T(z)T(w)\rangle_\tau^{\reg}}{\omega_1(z)}
\;=\;
\Res^{(2)}_{w=z}\!\bigl(\langle T(z)T(w)\rangle_\tau\bigr)
\,/\,\omega_1(z),
\end{equation}
where $\Res^{(2)}_{w=z}$ denotes the coefficient of
$(z-w)^{-2}\,dz\wedge d\bar z$ in the genus-$1$ Arakelov
propagator expansion. The identity is independent of the
insertion point~$z$ and of the choice of Arakelov
trivialisation, and is covariant under
$\tau\mapsto(a\tau+b)/(c\tau+d)$.
\end{conjecture}

\begin{remark}[Proof obligation]
The missing step is an Arakelov/Faltings or Bismut--Gillet--Soulé
derivation identifying the regularised torus stress-tensor two-point
coefficient with the Virasoro lane $c/2$. Stress-tensor two-point
functions do not recover the programme's full level-trace invariant
for Heisenberg or affine Kac--Moody without additional current/OPE
data.
\end{remark}

\begin{remark}[Per-family verification]
\label{rem:cardy-arakelov-families}
Equation~\eqref{eq:cardy-arakelov} is a Virasoro-lane test. For
Heisenberg and affine Kac--Moody, the full scalar invariant is
extracted from current/OPE trace data and the Sugawara shift, not from
the stress-tensor two-point function alone. For principal $\cW_N$, the
claimed equality with $c(H_N-1)$ is a separate family computation.
\end{remark}

\begin{remark}[Chiral Chern--Weil content]
\label{rem:cardy-is-chern-weil}
Equation~\eqref{eq:cardy-arakelov} is the proposed Virasoro-lane
chiral Chern--Weil form of Cardy's identity. On the classical side, a connection's
curvature pairs against an invariant polynomial to yield a
Chern class; on the chiral side, the Maurer--Cartan element
$\Theta^{\Eone}_\cA$ restricted to $\Sigma_1$ pairs against
the Arakelov form $\omega_1$ to yield the Virasoro scalar. The
full programme invariant still requires the family-specific
extraction rule.
\end{remark}

\subsection{The chiral Chern--Weil map}

As $[\Sigma_g]$ varies over $\Mbar_g$, the bar complex
forms a family of curved cochain complexes.  The fibrewise curvature
defines the following typed trace data.

\begin{theorem}[Chiral Chern--Weil; \ClaimStatusConditional]\label{thm:chiral-cw}
\textup{[Type signature: Open quadrant; modular-trace presentation;
levels $4\to5$; package
$H_D=(H_D^1,H_D^K,H_D^{\mathrm{tr}},H_D^{\mathrm{graph}})$.]}
For a cyclic chart algebra~$\cA$ and $g\geq2$, the native class is
$\operatorname{Obs}^{\mathrm{def}}_g(\cA)\in
H^2(\operatorname{Def}_g(\cA))$.  Under~$H_D^K$,
\begin{equation}\label{eq:obs-class}
\mathfrak O_g^K(\cA)=\kappa(\cA)\lambda_{-1}(\mathbb E_g),
\qquad
\operatorname{ch}_g(\mathfrak O_g^K(\cA))
=(-1)^g\kappa(\cA)\lambda_g
 \quad \end{equation}
with $\lambda_g=c_g(\mathbb E_g)$.  Under~$H_D^1$ the pointed
genus-one trace of the deformation class is $\kappa(\cA)\lambda_1$;
$H_D^{\mathrm{tr}}$ supplies its comparison with~$\mathfrak O_g^K$.
Under~$H_D^{\mathrm{graph}}$ the numerical output is
$F_g=\kappa\lambda_g^{\mathrm{FP}}+\delta F_g^{\mathrm{cross}}$.
\end{theorem}

The chiral Chern--Weil map:
\begin{equation}\label{eq:chiral-cw-map}
\underbrace{\Theta^{\Eone}_\cA}_{\text{MC element}}
\;\xrightarrow{\;q\;}
\;\underbrace{\Theta^{\mathrm{mod}}_\cA}_{\text{homotopy-coinvariant MC element}}
\;\xrightarrow{\;\text{genus-$g$ deformation component}\;}
\;\underbrace{\operatorname{Obs}^{\mathrm{def}}_g\in
H^2(\operatorname{Def}_g)}_{\substack{\text{native deformation class}\\
g\geq2}}
\;\xrightarrow{\;H_D^{\mathrm{tr}}\;}
\;\underbrace{\mathfrak O_g^K\in K^0(\Mbar_g)}_{\text{perfect object}}
\;\xrightarrow{\;\operatorname{ch}_g\;}
\;\underbrace{(-1)^g\kappa\lambda_g\in H^{2g}(\Mbar_g)}_{\text{Hodge character}}.
\end{equation}

\begin{remark}[Comparison with classical Chern--Weil]
\label{rem:comparison}
\begin{center}
\renewcommand{\arraystretch}{1.3}
\begin{tabular}{@{}p{4.2cm}p{5.5cm}@{}}
\toprule
\textbf{Classical} & \textbf{Chiral} \\
\midrule
Connection $\nabla$ & MC element $\Theta^{\Eone}_\cA$ \\
Curvature $F_\nabla \in \Omega^2(M;\fg)$
 & $K_\cA^{\mathrm{coll}}(z) \in \gEone$; $\dfib^{\,2} = \kappa\cdot\omega_g$ \\
Inv.\@ polynomial $P \in \Sym^k(\fg^*)^G$
 & Averaging map $\av\colon\gEone\to\gmod$ \\
$\CW(P) = P(F^k)$
 & $\tau_{\cA,2}^{\mathrm{res}}q_2(K_\cA^{\mathrm{coll}})=\kappa_{\mathrm{cl}}$ \\
Bianchi $d_\nabla F = 0$
 & MC: $d\Theta + \frac{1}{2}[\Theta,\Theta] = 0$ \\
Char.\@ class $[P(F^k)] \in H^{2k}(M)$
 & $\operatorname{ch}_g(\mathfrak O_g^K)=(-1)^g\kappa\lambda_g$ under $H_D^K$ \\
Independence of connection
 & $\kappa$ is a qi-invariant \\
$H^*(BG) = \Sym^*(\fg^*)^G$
 & Shadow tower $(\kappa,\mathfrak{C},\mathfrak{Q},\ldots)$ \\
\bottomrule
\end{tabular}
\end{center}
Three structural features have no classical analogue:
(i)~the base curve~$X$ itself varies over~$\Mbar_g$;
(ii)~the curvature factorizes as algebra $\times$ geometry;
(iii)~the characteristic classes are graded by
genus~$g$, not by degree of an invariant polynomial.
\end{remark}


% ================================================================
% Section 6: The shadow tower
% ================================================================
\section{The shadow tower}\label{sec:shadow}

The modular characteristic~$\kappa$ is the scalar coordinate
$\tau_{\cA,2}^{\mathrm{res}}q_2(K_\cA^{\mathrm{coll}})$ in the abelian
and scalar lanes.  The ordered MC element
$\Theta^{\Eone}_\cA$ has components at every degree:
the KZ associator~$r_3$ at degree~$3$, the quartic coherence
$r_4$ at degree~$4$, and so on. Each carries $\Eone$ operadic
indices; each has a $\Sigma_n$-coinvariant shadow. The
sequence of these shadows is the \emph{shadow tower}, and
its depth classifies chiral algebras into four types.

\subsection{The shadow algebra}

\begin{definition}[Shadow algebra]\label{def:shadow-algebra}
Let $\cA$ be a chiral algebra with modular cyclic deformation
complex $\Defcyc^{\mathrm{mod}}(\cA)$. The \emph{shadow
algebra} is
\begin{equation}\label{eq:shadow-algebra}
\Ash := H_\bullet\bigl(\Defcyc^{\mathrm{mod}}(\cA)\bigr),
\end{equation}
equipped with the descended Lie bracket from the convolution
bracket on $\Defcyc^{\mathrm{mod}}(\cA)$.
\end{definition}

\begin{proposition}\label{prop:shadow-bracket}
The convolution bracket has degree~$0$ and maps
degree~$r_1 \times$ degree~$r_2$ to degree~$r_1 + r_2 - 2$
\textup{(}by gluing two external legs\textup{)}. It descends
to a well-defined graded Lie bracket on~$\Ash$, making
$(\Ash, [-,-])$ a bigraded Lie algebra.
\end{proposition}

\begin{proof}
The bracket $[-,-]$ on $\Defcyc^{\mathrm{mod}}(\cA)$ is a
chain map: $d[f,g] = [df,g] + (-1)^{|f|}[f,dg]$. Cocycles
map to cocycles; the bracket of a cocycle with a coboundary is
a coboundary. The descended bracket inherits graded
antisymmetry and the Jacobi identity.
\end{proof}

The shadow algebra carries a bigrading by degree $r \ge 2$ and
genus $g \ge 0$:
\begin{equation}\label{eq:shadow-bigrading}
\Ash = \bigoplus_{r \ge 2}\;\bigoplus_{g \ge 0}\;
\Ash_{r,g}.
\end{equation}

\subsection{Shadow projections}

The individual shadow invariants are the bigraded components:
\begin{itemize}
\item $\kappa(\cA) = \Ash_{2,0}$: the \emph{modular
 characteristic} (degree~$2$, genus~$0$).
\item $\Delta_\cA = \Ash_{2,*}$: the \emph{spectral
 discriminant} (all genera at degree~$2$).
\item $\mathfrak{C}(\cA) = \Ash_{3,0}$: the \emph{cubic
 shadow} (degree~$3$, genus~$0$).
\item $\mathfrak{Q}(\cA) = \Ash_{4,0}$: the \emph{quartic
 shadow} (degree~$4$, genus~$0$), containing~$S_4$.
\item $\Sh_r(\cA) = \Ash_{r,0}$: the degree-$r$ shadow.
\end{itemize}

Each scalar shadow is obtained by applying the appropriate normalized
functional to the corresponding homotopy-coinvariant operator class:
\begin{equation}\label{eq:shadow-from-e1}
(\kappa,\;\mathfrak{C},\;\mathfrak{Q},\;\ldots)
= \bigl(
\tau_{\cA,2}^{\mathrm{res}}q_2(K_\cA^{\mathrm{coll}}),
\tau_{\cA,3}^{\mathrm{res}}q_3((\Theta_\cA^{\Eone})_{0,3}),
\tau_{\cA,4}^{\mathrm{res}}q_4((\Theta_\cA^{\Eone})_{0,4}),
\ldots\bigr).
\end{equation}

\subsection{What the higher-degree objects are}

\noindent
\emph{Degree~$2$}: the $\Eone$ datum is the collision kernel
$K_\cA^{\mathrm{coll}}$.  Its homotopy coinvariant is the operator
$q_2(K_\cA^{\mathrm{coll}})$, and the scalar coordinate is
$\tau_{\cA,2}^{\mathrm{res}}q_2(K_\cA^{\mathrm{coll}})$.  Under
$H_{\mathrm{CYBE}}^{\mathrm{rep}}$, its represented image is a
classical Yang--Baxter datum.

\medskip\noindent
\emph{Degree~$3$}: the $\Eone$ datum is the KZ associator
$r_3(z_1,z_2)$, a function on $\Conf_3(X)$ satisfying the
pentagon identity. Its image under~$\av$ is the cubic
shadow~$\mathfrak{C}$: the irreducible triple-collision
invariant.

\medskip\noindent
\emph{Degree~$4$}: the $\Eone$ datum is the quartic
$R$-matrix coherence~$r_4$, satisfying the quartic Stasheff
equation. Its image under~$\av$ contains the quartic contact
invariant~$S_4$.

\medskip\noindent
\emph{General degree~$r$}: the obstruction class at degree~$r+1$
is
\begin{equation}\label{eq:obstruction-class}
o_{r+1}(\cA)
= \bigl(D_\cA\Theta^{\le r}_\cA
+ \tfrac{1}{2}[\Theta^{\le r}_\cA,
 \Theta^{\le r}_\cA]\bigr)_{r+1}
\;\in\; H^2(\Ash_{r+1,0}).
\end{equation}
This is a Lie bracket of shadows at degrees~$\le r$: the
shadow algebra bracket connects different degree components and
drives the MC recursion.

\subsection{The discriminant and the G/L/C/M classification}

\begin{definition}\label{def:discriminant}
The \emph{discriminant} of a chiral algebra~$\cA$ is
\begin{equation}\label{eq:discriminant}
\Delta = 8\kappa\cdot S_4.
\end{equation}
\end{definition}

\begin{theorem}[Depth classification]\label{thm:depth-class}
For the standard verified families, the tuple
$(S_3,S_4,\text{termination})$ classifies the four coarse depth
types:
\begin{center}
\begin{tabular}{@{}lllll@{}}
\toprule
Class & Depth & Shadow criterion & Prototype & $\Delta$ \\
\midrule
$G$ \textup{(Gaussian)} & $2$ & $S_3=S_4=0$ &
 Heisenberg & $0$ \\
$L$ \textup{(Lie)} & $3$ & $S_3\ne0,\ S_4=0$ &
 Affine KM & $0$ \\
$C$ \textup{(Contact)} & $4$ &
 $S_3=0,\ S_4\ne0,\ S_{r\ge5}=0$ &
 $\beta\gamma$ & $\ne0$ generically \\
$M$ \textup{(Mixed)} & $\infty$ & non-terminating &
 $\Vir$, $\cW_N$ & $\neq 0$ \\
\bottomrule
\end{tabular}
\end{center}
\end{theorem}

\begin{remark}
The discriminant $\Delta=8\kappa S_4$ is downstream data. It detects
the quartic/contact channel when it is defined; it is not by itself an
archetype classifier.
\end{remark}

\begin{remark}[Analogy with classifying spaces]
\label{rem:shadow-bg}
The G/L/C/M classification is the chiral analogue of
classifying principal bundles by the type of their structure
group. The number of independent shadow invariants
(characteristic classes) determines the complexity of the
chiral classifying space. Class~$G$
(one invariant,~$\kappa$) is analogous to $BU(1)$ (one
class,~$c_1$). Class~$M$ (infinitely many invariants) is
analogous to $BO(\infty)$ (infinitely many Pontryagin classes).

A structural difference: classical Chern classes are
polynomials in a single curvature~$F$ (the degree-$2$
datum); chiral shadow invariants are averages of distinct
higher-degree operations~$r_n$. Each shadow captures
genuinely new collision structure not determined by lower
degrees, unless $\Delta = 0$ and the tower terminates.
\end{remark}

\begin{remark}[Shadow depth and Koszulness]
\label{rem:depth-koszul}
Shadow depth measures how many shadow invariants are
non-zero. Koszulness (bar cohomology concentrated in
degree~$1$) is a separate property. All standard families are
Koszul; the shadow depths differ. SC-formality
($m_k^{SC} = 0$ for $k \ge 3$) holds if and only if the
algebra is class~$G$.
\end{remark}


% ================================================================
% Section 7: Multi-weight
% ================================================================
\section{Multi-weight structure and cross-channel
corrections}\label{sec:multi-weight}

The package~$H_D^1$ gives the pointed genus-one trace
$\kappa\lambda_1$.  The package~$H_D^K$ gives the perfect object and
its signed Hodge character at every genus in its scope.  The package
$H_D^{\mathrm{graph}}$ gives the full numerical functional, including
the mixed-weight cross-channel summand.  The first witness is $\cW_3$
at genus~$2$,
where $\delta F_2 = (c{+}204)/(16c)$. This section derives
these results from the graph-sum structure of the genus
expansion.

\subsection{Per-channel decomposition}

Let $\cA$ be a modular Koszul chiral algebra with strong
generators $\varphi_1,\ldots,\varphi_r$ of conformal weights
$h_1,\ldots,h_r$. Each generator has a per-channel modular
characteristic~$\kappa_i$: the contribution from the
$(\varphi_i,\varphi_i)$ self-OPE channel. The total is
$\kappa(\cA) = \sum_{i=1}^r\kappa_i$.

The genus-$g$ free energy is computed by a graph sum over
stable graphs~$\Gamma$ of genus~$g$:
\begin{equation}\label{eq:graph-sum}
F_g(\cA) = \sum_\Gamma |\Aut(\Gamma)|^{-1}
\sum_{\sigma\colon E(\Gamma) \to \{1,\ldots,r\}}
A_\Gamma(\sigma,\cA),
\end{equation}
where $\sigma$ assigns a channel to each edge~$e$
of~$\Gamma$.

\begin{definition}\label{def:diagonal-mixed}
A channel assignment~$\sigma$ is \emph{diagonal} if
$\sigma(e) = i$ for all~$e$ (all edges carry the same channel).
It is \emph{mixed} if at least two edges carry different
channels.
\end{definition}

\subsection{The decomposition theorem}

\begin{theorem}[Conditional multi-weight genus expansion]
\label{thm:multi-weight}
Assume the diagonal genus-$g$ Faber--Pandharipande comparison and the
mixed-channel boundary graph expansion of Volume I. With the notation
above:
\begin{enumerate}[\upshape(i)]
\item \emph{Per-channel universality.}
 The diagonal contribution satisfies
 \begin{equation}\label{eq:diag}
 F_g^{\mathrm{diag}}(\cA) = \kappa(\cA)\cdot
 \lambda_g^{\mathrm{FP}},
 \end{equation}
 where $\lambda_g^{\mathrm{FP}}
 = \frac{2^{2g-1}-1}{2^{2g-1}}\cdot
 \frac{|B_{2g}|}{(2g)!}$ is the Faber--Pandharipande
 number ($B_{2g}$ the Bernoulli number), for all
 $g \ge 1$ in the diagonal sector.

\item \emph{Cross-channel decomposition.}
 The full free energy decomposes as
 \begin{equation}\label{eq:decomposition}
 F_g(\cA) = \kappa(\cA)\cdot\lambda_g^{\mathrm{FP}}
 + \delta F_g^{\mathrm{cross}}(\cA),
 \end{equation}
 where $\delta F_g^{\mathrm{cross}}$ is a graph sum over
 mixed-channel boundary graphs of~$\Mbar_{g,0}$.

\item \emph{Genus-$1$ universality.}
 $\delta F_1^{\mathrm{cross}} = 0$ for all~$\cA$.

\item \emph{Uniform-weight universality.}
 If $h_1 = \cdots = h_r$, then
 $\delta F_g^{\mathrm{cross}} = 0$ at all $g \ge 1$.

\item \emph{$R$-matrix scope.}
 No higher-genus mixed-channel $R$-independence is asserted here.
 The $R$-independence statement is restricted to diagonal terms and
 to genus~$1$.

\item \emph{Genus-$2$ explicit formula.}
 At genus~$2$:
 \begin{equation}\label{eq:genus2-cross}
 \delta F_2^{\mathrm{cross}}(\cA)
 = \delta F_2^{\mathrm{sun}}(\cA)
 + \delta F_2^{\theta}(\cA)
 + \delta F_2^{\mathrm{lt}}(\cA)
 + \delta F_2^{\mathrm{bb}}(\cA),
 \end{equation}
 summing over four contributing boundary strata among the
 seven stable graphs of~$\Mbar_{2,0}$.
\end{enumerate}
\end{theorem}

\begin{proof}
(i)~On a diagonal assignment $\sigma(e) = i$ for all~$e$,
the amplitude factors as $A_\Gamma(\sigma_i,\cA)
= A_\Gamma(\cA_i)$ for the rank-$1$ algebra~$\cA_i$ with
curvature~$\kappa_i$. Each rank-$1$ component has a single
generator and hence uniform weight. Summing:
\[
F_g^{\mathrm{diag}}(\cA) = \sum_{i=1}^r F_g(\cA_i)
= \sum_{i=1}^r \kappa_i\cdot\lambda_g^{\mathrm{FP}}
= \kappa(\cA)\cdot\lambda_g^{\mathrm{FP}}.
\]

\smallskip\noindent
(iii)~At genus~$1$, the unique boundary stratum has a single
edge. A single edge admits only $r$ channel assignments,
each pure. No mixed-channel assignments exist on a
single-edge graph.
\end{proof}

\subsection{Why cross-channel corrections arise at genus g >= 2}\label{subsec:why-cross}

Three ingredients are needed for a non-zero
$\delta F_g^{\mathrm{cross}}$:
\begin{enumerate}[(a)]
\item \emph{Multiple edges on a single graph}, so that
 different edges can carry different channels. This
 requires $g \ge 2$, since all genus-$1$ contributing graphs
 have a single edge.

\item \emph{Distinct weights} on different channels
 ($h_i \neq h_j$), so that the propagators along edges
 carry different channel weights in the universal propagator.

\item \emph{Non-trivial OPE mixing}: the vertices carry
 mixed-channel OPE structure constants not controlled
 by~$\kappa$ alone.
\end{enumerate}

At genus~$1$, ingredient~(a) fails: the single-edge graph
cannot carry a mixed assignment.
At uniform weight, ingredient~(b) fails: all propagators are
identical regardless of channel.
For free fields, ingredient~(c) fails: $m_k = 0$ for
$k \ge 3$ at genus~$0$, so mixed-channel vertices carry no
structure constants beyond the two-point function.

\subsection{Explicit computation: $\cW_3$ at genus~$2$}
\label{subsec:w3-genus2}

$\cW_3$ has two generators: $T$ (weight~$2$) and~$W$
(weight~$3$). Per-channel modular characteristics:
$\kappa_T = c/2$, $\kappa_W = c/3$. Total:
$\kappa = c/2 + c/3 = 5c/6$.

The four contributing boundary strata at genus~$2$ give:
\begin{equation}\label{eq:w3-cross}
\delta F_2(\cW_3) = \frac{3}{c} + \frac{9}{2c}
+ \frac{1}{16} + \frac{21}{4c}
= \frac{c + 204}{16c}.
\end{equation}

This is non-zero for $c\ne0,-204$; for positive central charge it is
non-zero.

\begin{itemize}
\item At $c = 1$: $\delta F_2 = 205/16$.
\item At large~$c$: $\delta F_2 \to 1/16$ (the leading
 correction is a constant, independent of~$c$).
\item At $c = -204$: $\delta F_2 = 0$ (an isolated zero,
 not physically relevant).
\end{itemize}

The four terms arise from distinct boundary mechanisms:
\begin{itemize}
\item $3/c$: sunset graph with mixed-weight vertex carrying
 the $T$--$W$ OPE coefficient.
\item $9/(2c)$: theta graph with three edges at mixed weight
 assignments $(2,2,3)$, $(2,3,3)$, and permutations.
\item $1/16$: a topological contribution from the difference
 in Euler characteristics of the weight-$2$ and weight-$3$
 determinant line bundles over~$\Mbar_2$, independent of~$c$.
\item $21/(4c)$: bridge-loop graph connecting a genus-$1$
 vertex (single-edge loop at weight~$h_i$) to a genus-$0$
 vertex through an edge at weight~$h_j \neq h_i$.
\end{itemize}

\begin{corollary}\label{cor:scalar-fails}
The scalar formula $F_g(\cA) = \kappa(\cA)\cdot
\lambda_g^{\mathrm{FP}}$ fails at $g = 2$ for~$\cW_3$.
This resolves the multi-generator universality question
in the negative: the scalar formula is not universal for
interacting multi-weight algebras.
\end{corollary}

\subsection{Free-field exactness}

\begin{proposition}[Free-field exactness]
\label{prop:free-field-exact}
Let $\cA$ be a free-field chiral algebra
\textup{(}$m_k = 0$ for all $k \ge 3$ at genus~$0$\textup{)}.
Then $\delta F_g^{\mathrm{cross}}(\cA) = 0$ for all
$g \ge 1$. The scalar formula
$F_g = \kappa\cdot\lambda_g^{\mathrm{FP}}$ is exact at all
genera.
\end{proposition}

Three independent mechanisms force this vanishing:
\begin{enumerate}[(1)]
\item The genus-$0$ OPE is purely quadratic, so mixed-channel
 vertices carry no structure constants beyond the two-point
 function.
\item The Givental $R$-matrix is trivial for free fields,
 eliminating $R$-corrections at genus-$0$ vertices.
\item The propagator factorization identity for free fields
 collapses the cross-channel graph sum to the diagonal.
\end{enumerate}


% ================================================================
% Section 8: The chiral classifying space
% ================================================================
\section{The chiral classifying space}\label{sec:classifying}

The preceding sections constructed each ingredient of chiral
Chern--Weil theory individually: the collision kernel (the binary
connection datum), the homotopy-coinvariant projection, the normalized
trace functional (the scalar invariant), the curvature formula, the shadow tower (the
characteristic ring), and the multi-weight correction.
The forced question: is there a single geometric object that
unifies all of these, the way the classifying space~$BG$
unifies classical Chern--Weil theory?

\subsection{$B(\cA)$ as classifying object}

The ordered bar complex $B^{\mathrm{ord}}(\cA)$ on $\Ran(X)$,
equipped with the deconcatenation coproduct, is a factorization
coalgebra. In the deformation-theoretic picture, it is the
chiral analogue of the simplicial bar construction
$[n] \mapsto G^n$ whose geometric realization is the classifying
space~$BG$.

The universal bundle is the twisted tensor product
\begin{equation}\label{eq:universal-bundle}
\cE^{\mathrm{univ}}(\cA) = \cA \otimes^\tau B(\cA),
\end{equation}
with differential combining the internal differential, the bar
differential, and the twisting morphism
$\tau\colon B(\cA) \to \cA$. Its acyclicity,
$\Omega(B(\cA)) \xrightarrow{\sim} \cA$ (bar-cobar
quasi-isomorphism on the Koszul locus), is the chiral analogue
of the contractibility of~$EG$.

\subsection{The universal connection}

The MC element $\Theta_\cA \in \MC(\gmod)$ is the
\emph{universal connection} on the chiral classifying space.
The MC equation
\begin{equation}\label{eq:mc-eq}
d\Theta_\cA + \tfrac{1}{2}[\Theta_\cA,\Theta_\cA] = 0
\end{equation}
is the flatness/integrability condition. At genus~$0$, the
universal connection is flat ($\dfib^{\,2} = 0$); the
universal bundle is trivializable over~$\cM_{0,n}$.

At genus~$g\geq2$, the fibrewise curvature is
$\dfib^{\,2}=\kappa\omega_g$.  The modular trace assigns the native
class $\operatorname{Obs}^{\mathrm{def}}_g$, the perfect object
$\mathfrak O_g^K=\kappa\lambda_{-1}(\mathbb E_g)$ under~$H_D^K$,
and the numerical graph functional under~$H_D^{\mathrm{graph}}$.
At genus one, $H_D^1$ begins with the pointed class
$\operatorname{Obs}^{\mathrm{def}}_{1,1}$.

\subsection{The total system over moduli}

The total system is a family
\begin{equation}\label{eq:total-system}
\cE_{\mathrm{total}}
\;\longrightarrow\;
B^{\Sigma}(\cA)\big|_{\Sym^\bullet(\mathcal{C}_g/\Mbar_g)}
\;\longrightarrow\;
\Mbar_g,
\end{equation}
where $\mathcal{C}_g \to \Mbar_g$ is the universal curve, and
$B^{\Sigma}(\cA)$ is the $\Sigma_n$-coinvariant (symmetric) bar
complex over the symmetric product
$\Sym^n(X) = \Conf_n(X)/\Sigma_n$.

The ordered bar $B^{\mathrm{ord}}(\cA)$ on $\Conf_n(X)$ carries
$K_\cA^{\mathrm{coll}}(z)$ and, on represented comparison surfaces,
the corresponding $R$-operator. The symmetric bar
$B^{\Sigma}(\cA)$ on $\Sym^n(X)$ retains the homotopy-coinvariant
operator. The normalized trace gives $\kappa$ on scalar families and
$\kappa_{\mathrm{dp}}$ on non-abelian affine families, where the
Sugawara contribution completes $\kappa$.

As $[X]$ moves in~$\Mbar_g$:
\begin{enumerate}[(i)]
\item The bar differential depends on~$X$ through the prime-form
 propagator $d\log E(z,w)$ (weight~$1$).
\item At genus~$0$: flat ($\dfib^{\,2} = 0$), universal bundle
 trivializable.
\item At genus~$g \ge 1$: curvature
 $\dfib^{\,2}=\kappa\omega_g$; $H_D^K$ gives
 $\mathfrak O_g^K=\kappa\lambda_{-1}(\mathbb E_g)$.
\item At multi-weight, genus $g \ge 2$: $H_D^{\mathrm{graph}}$ gives
 $F_g=\kappa\lambda_g^{\mathrm{FP}}+\delta F_g^{\mathrm{cross}}$.
\end{enumerate}

\subsection{Crystallization}

Classical Chern--Weil theory has three ingredients: a
connection~$\nabla$, an invariant polynomial~$P$, and the
resulting characteristic class~$[P(F^k)]$. Chiral
Chern--Weil theory has the same three: the $\Eone$ MC
element~$\Theta^{\Eone}_\cA$ (whose binary component is
$K_\cA^{\mathrm{coll}}$), the homotopy-coinvariant map
$q_n$, the normalized functional $\tau_{\cA,n}^{\mathrm{res}}$,
and, for $g\geq2$, the Theorem~D objects
$\operatorname{Obs}^{\mathrm{def}}_g$, $\mathfrak O_g^K$, its signed
Hodge character, and~$F_g$.
What is new is the variation over moduli: the curvature
$\dfib^{\,2} = \kappa\cdot\omega_g$ factorizes as
algebra~$\times$ geometry, the shadow tower
$(\kappa,\mathfrak{C},\mathfrak{Q},\ldots)$ replaces the
finite-rank invariant ring of a compact group with an
degree-graded sequence whose depth classifies chiral algebras
(G/L/C/M), and the scalar formula fails at genus~$\ge 2$ for
multi-weight algebras with explicit correction
$\delta F_2(\cW_3) = (c{+}204)/(16c)$.


% ================================================================
% Koszul conductor and scalar complementarity
% ================================================================
\section{Koszul conductors and scalar complementarity}
\label{sec:conductor-scalar}

The chiral Chern--Weil output~$\kappa(\cA)$ also organises
Koszul-complementary pairs. Two distinct invariants live on a
Koszul-complementary pair:
\begin{equation}\label{eq:universal-conductor}
K_c(\cA) \;:=\; c(\cA) + c(\cA^!)
\end{equation}
and the scalar $\kappa$-complementarity
$K^{\kappa}(\cA) := \kappa(\cA) + \kappa(\cA^{!})$. There is no
single ratio converting $K_c$ to $K^\kappa$: the affine
line has $K_c=2\dim\fg$ but $K^\kappa=0$, while Virasoro,
principal $\cW_N$, and Bershadsky--Polyakov use family-specific
normalisations.
	Trinity values: $K_c = 0$ for class~G and abelian Kac--Moody,
	$K_c = 2\dim(\fg)$ for non-abelian $V_k(\fg)$ while
	$K^\kappa=0$ on the affine line; $K_c = 26$ for Virasoro
	(so $K^{\kappa} = 13$ at self-dual
	$c = 13$, with $\varrho_{\mathrm{Vir}} = 1/2$); $K_c = 50$ for the
Bershadsky--Polyakov algebra $\mathrm{BP}$ (the polynomial
identity $\mathrm{K}_{\mathrm{BP}}(k) \equiv 50$ as a rational
function in the level, in the standard Fehily--Kawasetsu--Ridout
convention, with the canonical
fixed point at $k = -3 = -h^\vee(\mathfrak{sl}_3)$, where
conditionally
$K^{\kappa}_{\mathrm{BP}} = 25/3$ and $\varrho_{\mathrm{BP}} = 1/6$);
$K = 100$ for principal $\cW_3$ (so $K^{\kappa} = 250/3$ with
$\varrho_{\cW_3} = 5/6 = H_3 - 1$). The tabulated central charges of
the bosonic $\beta\gamma$ and
fermionic $bc$ ghost systems provide the universal base:
$c^{\beta\gamma}(\lambda) = 2(6\lambda^2 - 6\lambda + 1)$ and
$c^{bc}(\lambda) = 1 - 3(2\lambda - 1)^2$, whose sum vanishes
identically.

\begin{proposition}[Verified conductor lanes]
\label{thm:koszul-conductor}
For the verified standard families, central-charge conductors and
scalar modular conductors are family-specific quantities. There is no
single BRST ghost object canonically attached to every family in the
table.
\end{proposition}

\begin{proof}
For affine $V_k(\fg)$, the central-charge conductor and the scalar
modular conductor are different quantities:
\[
K_c(V_k(\fg))
:=c(V_k(\fg))+c(V_{-k-2h^\vee}(\fg))
=\frac{k\dim\fg}{k+h^\vee}
+\frac{(-k-2h^\vee)\dim\fg}{-k-h^\vee}
=2\dim\fg,
\]
whereas
\[
K^\kappa(V_k(\fg))
:=\kappa(V_k(\fg))+\kappa(V_{-k-2h^\vee}(\fg))
=\frac{\dim\fg}{2h^\vee}\bigl((k+h^\vee)+(-k-h^\vee)\bigr)
=0.
\]
The ghost calculation records the central-charge lane; the
vanishing scalar lane is the affine complementarity used by the
shadow tower. For $\Vir_c$: the $(b,c)$ ghost system has
$\lambda = 2$ giving $c^{bc}(2) = -26$; complementarity
$\kappa(\Vir_c) + \kappa(\Vir_{26-c}) = c/2 + (26-c)/2 = 13 = -c_{\mathrm{ghost}}/2$
after the universal factor. For $\cW_N$: the reduction from
$\fsl_N$ contributes a tower of ghosts at weights
$\{2,3,\ldots,N\}$; summation reproduces the tabulated $K$-value.
For $\mathrm{BP}$, the
polynomial identity $K_{\mathrm{BP}}(k) \equiv 50$ in the standard
Fehily--Kawasetsu--Ridout convention matches the central-charge
conductor $K_c^{\mathrm{BP}}=50$ at every~$k$; the conditional scalar
conductor on the ratio lane is
$K^\kappa_{\mathrm{BP}}=50/6=25/3$, and the fixed point
$k = -h^\vee(\mathfrak{sl}_3) = -3$ is a pole with
principal-value half $25/6$.
\end{proof}

\begin{proposition}[Scalar complementarity and moment residues at genus zero]
\label{thm:moment-genus-zero}
On the verified diagonal standard families, the genus-zero extraction
of the scalar invariant uses the moment-residue pairing
\begin{equation}\label{eq:moment-g0}
\kappa(\cA)
\;=\;
\sum_i \Res_{z=0}\!\bigl(z^{2h_i-2}K^{\mathrm{vac}}_{ii}(z)\,dz\bigr)
\;+\;\kappa_{\mathrm{sp}},
\end{equation}
with the affine Sugawara shift $\kappa_{\mathrm{sp}}$ added
separately. For a Koszul pair, the scalar conductor is then
$K^\kappa(\cA)=\kappa(\cA)+\kappa(\cA^!)$.
\end{proposition}

\begin{remark}[Epistemic tags and independent-verification conventions]
\label{rem:conductor-epistemic}
Every collision kernel in this section carries its family parameter
explicitly: the Kac--Moody trace-form kernel is
$K_{k,\fg}^{\mathrm{coll}}(z)=k\,\Omega_{\mathrm{tr}}/z$, and the
Virasoro kernel
$K_{\Vir_c}^{\mathrm{coll}}(z)=(c/2)/z^3+2T/z$ has the canonical
cubic pole. Every $\kappa$ value is taken from the landscape census.
The scalar conductor values $\{0, 13, 25/3, 250/3\}$ are verified
against two disjoint sources under the independent verification
protocol: direct computation of $\kappa + \kappa^!$ and ghost
tabulation after conversion from $K_c$ to $K^\kappa$ where a ghost
description is available.
\end{remark}

\begin{remark}[Cross-volume bridge]
\label{rem:conductor-cross-volume}
The central and scalar conductor lanes are separate Vol~I
manifestations of the trace formalism. In Vol~III, the scalar
normalization reappears
on Borcherds--Kac--Moody algebras as
$\kappa_{\mathrm{BKM}}(\Phi_N) = c_N(0)/2$ with $N \in \{1,2,3,4,6\}$
the level of the Fricke-type paramodular form. A numerical equality
with a Vol~I scalar conductor is not asserted without specifying the
Borcherds input $\Phi_N$, the normalization, and the Vol~I family
being compared.
\end{remark}


% ================================================================
% References (placeholder)
% ================================================================
\begin{thebibliography}{99}

\bibitem{BeiDr}
A.~Beilinson, V.~Drinfeld.
\emph{Chiral Algebras}.
AMS Colloquium Publications~\textbf{51}, 2004.

\bibitem{CG}
K.~Costello, O.~Gwilliam.
\emph{Factorization Algebras in Quantum Field Theory},
Vols.~1--2.
Cambridge University Press, 2017, 2021.

\bibitem{ChSe}
C.~Chevalley, S.~Eilenberg.
Cohomology theory of Lie groups and Lie algebras.
\emph{Trans.~AMS}~\textbf{63} (1948), 85--124.

\bibitem{CS}
S.~S.~Chern, J.~Simons.
Characteristic forms and geometric invariants.
\emph{Ann.~Math.}~\textbf{99} (1974), 48--69.

\bibitem{Drinfeld}
V.~G.~Drinfeld.
Quasi-Hopf algebras.
\emph{Leningrad Math.~J.}~\textbf{1} (1990), 1419--1457.

\bibitem{FeiginFrenkel}
B.~Feigin, E.~Frenkel.
Affine Kac--Moody algebras at the critical level and
Gel'fand--Dikii algebras.
\emph{Int.~J.~Mod.~Phys.~A}~\textbf{7}, Suppl.~1A (1992),
197--215.

\bibitem{FLuk}
V.~A.~Fateev, S.~L.~Lukyanov.
The models of two-dimensional conformal quantum field theory
with $\mathbb{Z}_n$ symmetry.
\emph{Int.~J.~Mod.~Phys.~A}~\textbf{3} (1988), 507--520.

\bibitem{FBZ}
E.~Frenkel, D.~Ben-Zvi.
\emph{Vertex Algebras and Algebraic Curves}.
Math.~Surveys Monographs~\textbf{88}, AMS, 2004.

\bibitem{Kac}
V.~Kac.
\emph{Infinite-Dimensional Lie Algebras}.
Cambridge University Press, 3rd ed., 1990.

\bibitem{KV}
M.~Kontsevich, Y.~Soibelman.
Deformations of algebras over operads and the Deligne conjecture.
\emph{Conf.~Moshe Flato 2000, Vol.~1},
Math.~Phys.~Stud.~\textbf{21}, Kluwer, 2000, 255--307.

\bibitem{LV}
J.-L.~Loday, B.~Vallette.
\emph{Algebraic Operads}.
Grundlehren~\textbf{346}, Springer, 2012.

\bibitem{LorgatI}
R.~Lorgat.
\emph{Modular Koszul Duality}.
Perimeter Institute, 2026, in preparation.

\bibitem{LorgatShadow}
R.~Lorgat.
Modular Koszul duality and the shadow tower.
Perimeter Institute, 2026.

\bibitem{LorgatVirR}
R.~Lorgat.
The Virasoro $r$-matrix.
Perimeter Institute, 2026.

\bibitem{LorgatSevenFaces}
R.~Lorgat.
Seven faces of the $r$-matrix.
Perimeter Institute, 2026.

\bibitem{MilnorStasheff}
J.~Milnor, J.~Stasheff.
\emph{Characteristic Classes}.
Annals of Mathematics Studies~\textbf{76}, Princeton, 1974.

\bibitem{Stasheff}
J.~Stasheff.
Homotopy associativity of $H$-spaces, I, II.
\emph{Trans.~AMS}~\textbf{108} (1963), 275--312.

\bibitem{Weil}
A.~Weil.
\emph{G\'eom\'etrie diff\'erentielle des espaces fibr\'es}.
Unpublished manuscript, 1949.

\bibitem{Zam}
A.~B.~Zamolodchikov.
Infinite additional symmetries in two-dimensional conformal
quantum field theory.
\emph{Theor.~Math.~Phys.}~\textbf{65} (1985), 1205--1213.

\end{thebibliography}

\end{document}
